% ============================================================================
% sec17.tex — §17 Laboratory & Measurement Handbook
% Target: 12 pages | Figures: fig:17:thermometer-response, fig:17:probe-depth-bias
% ============================================================================
\binderchapter{17}{Laboratory \& Measurement Handbook (Thermometer Physics, Probe Time Constant $5\tau$, Immersion Bias, Calibration \& GUM Uncertainty Propagation)}{Laboratory \& Measurement Handbook}{sec:17}
\scopebanner{REQUIRED}[2027 Rules 4.f, 5.I.b, 5.II.c,g \& FAQs 09/30--10/06/2026 on stirring and probe depth (All Tournament Levels)]

\idxtopic{thermometry!overview}\idxtopic{laboratory measurement handbook}\idxdev{thermometer protocol (Rule 4.f)}
Accurate thermal metrology links Part~I written laboratory analysis (\sref{sec:18}) to Part~II 2027 Heat-Lamp Collector temperature prediction (\sref{sec:15}--\sref{sec:16}) \src{S01,S05}. Under 2027 Rule~\texttt{4.f}, the Event Supervisor (ES) provides a non-mercury thermometer reading to $\SI{0.1}{\celsius}$ (\emph{liquid-in-glass spirit, electronic digital probe, or non-contact infrared}) and uses the \textbf{same thermometer} for a team's starting ($T_{\mathrm{start}}$) and final ($T_f$) water readings \src{S01,S02}. This chapter covers the sensor physics of all five thermometer families, dynamic thermal lag ($\tau_p$ and the $5\tau_p$ settling rule), top-layer immersion and radial stratification bias under the FAQs on stirring and probe depth (09/30--10/06/2026), fixed-point ice/steam calibration, JCGM GUM / NIST TN~1297 uncertainty propagation \src{S21}, and step-by-step classroom-safe calorimetry and heat-lamp collector validation protocols.

% ============================================================================
\section{Physics \& Specifications of the Five Thermometer Technologies (\texttt{Rule 4.f})}\label{sec:17.1}
\idxtopic{thermometer technologies}\idxtopic{Rule 4.f!thermometer types}

Every thermometer transduces temperature $T$ into a measurable thermometric property $X(T)$ (liquid column length $L$, electrical resistance $R$, thermoelectric voltage $\Delta V$, or mid-infrared radiant exitance $M_{\lambda}$) \src{S13,S22,S33}. Table~\ref{tab:17:sensors} summarizes the physical laws, operating windows, dynamic time constants $\tau_p$, active sensing zones, and 2027 competition diagnostics across all five sensor technologies permitted under Rule~\texttt{4.f} \src{S01,S05,S22}.

\begin{table}[h!]
\centering
\captionof{table}{Comparison of the Five Thermometer Technologies Used in Science Olympiad Thermodynamics (\texttt{Rule 4.f}) \src{S01,S05,S12,S22}.}\label{tab:17:sensors}
\footnotesize
\setlength{\tabcolsep}{3pt}
\renewcommand{\arraystretch}{1.15}
\begin{tabularx}{\linewidth}{@{}>{\raggedright\arraybackslash}p{0.17\linewidth} >{\raggedright\arraybackslash}p{0.27\linewidth} C{0.10\linewidth} C{0.10\linewidth} C{0.09\linewidth} C{0.10\linewidth} >{\raggedright\arraybackslash}X@{}}
\toprule
\hdr{Sensor Family} & \hdr{Governing Thermometric Equation} & \hdr{Range $[\si{\celsius}]$} & \hdr{Typical Accuracy} & \hdr{$\tau_p$ in Water ($5\tau_p$)} & \hdr{Sensing Zone $[\si{\milli\meter}]$} & \hdr{2027 Top-Layer Stratification \& Exam Role} \\
\midrule
\textbf{1.~Liquid-in-Glass (Spirit / Galinstan)} \idxtopic{liquid-in-glass thermometer}\idxtopic{spirit thermometer} &
$L(T) = L_0 + \dfrac{V_0}{A_c}(\beta_{\mathrm{liq}} - 3\alpha_g)(T - T_0)$ &
$-20$ to $+110$ (spirit) &
$\pm\SIrange{0.5}{1.0}{\celsius}$ &
$\SIrange{6}{12}{\second}$ ($\SIrange{30}{60}{\second}$) &
$15$--$25$ bulb ($50$--$76$ imm.) &
Long bulb averages across top \& middle layers; stem-emergence \& meniscus parallax bias. Mercury banned (\texttt{Rule 4.f}). \\
\rowcolor{rowalt}
\textbf{2.~NTC Thermistor Digital Probe} \idxtopic{thermistor (NTC)}\idxformula{Steinhart-Hart equation} &
$\dfrac{1}{T} = A + B\ln R + C(\ln R)^3$; \newline $R(T) = R_0 e^{\beta(1/T - 1/T_0)}$ &
$-50$ to $+150$ &
$\pm\SIrange{0.1}{0.2}{\celsius}$ &
$\SIrange{0.8}{3.5}{\second}$ ($\SIrange{4}{18}{\second}$) &
$4$--$10$ (bead at tip) &
\textbf{Most common ES probe.} High sensitivity ($\alpha_R \approx \SI{-4}{\percent\per\kelvin}$) \& small tip sample top $\SIrange{5}{12}{\milli\meter}$ layer (FAQ 10/06/2026 on probe depth). \\
\textbf{3.~Platinum RTD (\texttt{Pt100} / \texttt{Pt1000})} \idxtopic{RTD (Pt100/Pt1000)}\idxformula{Callendar-Van Dusen equation} &
$R(t) = R_0(1 + A t + B t^2)$ ($t \ge \SI{0}{\celsius}$, $R_0 = \SI{100.00}{\ohm}$) &
$-200$ to $+600$ &
Class A: $\pm(0.15 + 0.002|t|)$ &
$\SIrange{3.0}{8.0}{\second}$ ($\SIrange{15}{40}{\second}$) &
$15$--$30$ (wire/film element) &
Reference metrology standard (IEC~60751). Longer sheath requires $\ge\SI{25}{\milli\meter}$ depth; use 3/4-wire bridge for lead resistance. \\
\rowcolor{rowalt}
\textbf{4.~Thermocouple (\texttt{Type K} / \texttt{Type T})} \idxtopic{thermocouple (Type K, T)}\idxtopic{Seebeck effect} &
$\Delta V = \int_{T_{\mathrm{ref}}}^{T} S_{AB}(T')\,dT' \approx S_{AB}(T - T_{\mathrm{ref}})$ \src{S22} &
$-200$ to $+1260$ (K); $+350$ (T) &
$\pm\SI{1.1}{\celsius}$ (K); $\pm\SI{0.5}{\celsius}$ (T) &
$\SIrange{0.05}{0.3}{\second}$ bare ($\SIrange{1.5}{3}{\second}$ sheath) &
$1$--$5$ (point weld) &
Fastest point sensor ($5\tau_p < \SI{1.5}{\second}$ bare bead) with negligible heat-capacity drain; requires cold-junction compensation (CJC). \\
\textbf{5.~Non-Contact Infrared (IR) Gun} \idxtopic{infrared (IR) thermometer}\idxdev{IR thermometer emissivity trap} &
$T_{\mathrm{read}} = \left[T_{\mathrm{amb}}^4 + \dfrac{\varepsilon_{\mathrm{true}}}{\varepsilon_{\mathrm{set}}}(T_{\mathrm{obj}}^4 - T_{\mathrm{amb}}^4)\right]^{1/4}$ &
$-30$ to $+500$ &
$\pm\SIrange{1.5}{2.0}{\celsius}$ &
$\SIrange{0.15}{0.5}{\second}$ ($\SIrange{1.0}{2.5}{\second}$) &
$0.015$ (top $\SI{15}{\micro\meter}$ skin) &
Water $\varepsilon_w = 0.96$ matches $\varepsilon_{\mathrm{set}} = 0.95$. \textbf{Never shoot through plastic wrap/cup wall or at shiny Al foil ($\varepsilon=0.04$)!} \\
\bottomrule
\end{tabularx}
\end{table}

\subsection{1. Non-Mercury Liquid-in-Glass Thermometers (Spirit \& Galinstan)}
\idxtopic{liquid-in-glass thermometer!apparent expansion}\idxformula{apparent volumetric expansion beta_app = beta_liq - 3*alpha_glass}
Because mercury ($\mathrm{Hg}$) is prohibited for safety under Rule~\texttt{4.f} \src{S01}, competition liquid-in-glass thermometers use dyed organic spirits (ethanol, $\beta_{\mathrm{liq}} = \SI{1.12e-3}{\per\kelvin}$; toluene, $\beta_{\mathrm{liq}} = \SI{1.08e-3}{\per\kelvin}$; or kerosene, $\beta_{\mathrm{liq}} = \SI{0.99e-3}{\per\kelvin}$) or non-toxic liquid metal alloy Galinstan ($\mathrm{Ga\text{--}In\text{--}Sn}$, $\beta_{\mathrm{liq}} \approx \SI{1.83e-4}{\per\kelvin}$) enclosed in a thin borosilicate glass bulb ($\alpha_g = \SI{3.3e-6}{\per\kelvin}$, volumetric $\beta_g = 3\alpha_g = \SI{9.9e-6}{\per\kelvin}$) connected to a uniform capillary bore of diameter $d_c$ and cross-sectional area $A_c = \pi d_c^2/4$ \src{S13}.
\begin{itemize}
  \item \textbf{Differential Volumetric Expansion Law:} When the bulb of initial internal volume $V_0$ warms by $\Delta T = T - T_0$, the liquid volume expands by $\Delta V_{\mathrm{liq}} = V_0 \beta_{\mathrm{liq}}\Delta T$ while the glass bulb cavity expands by $\Delta V_{\mathrm{glass}} = V_0 (3\alpha_g)\Delta T$. Only the \emph{apparent} differential volume $\Delta V_{\mathrm{app}} = V_0(\beta_{\mathrm{liq}} - 3\alpha_g)\Delta T = V_0\beta_{\mathrm{app}}\Delta T$ is forced up the capillary stem of area $A_c$:
  \begin{equation}\label{eq:17:lig}
    \Delta L = L(T) - L_0 = \frac{V_0}{A_c}(\beta_{\mathrm{liq}} - 3\alpha_g)(T - T_0) = \frac{4 V_0 \beta_{\mathrm{app}}}{\pi d_c^2}(T - T_0) \implies S_L = \frac{\Delta L}{\Delta T} = \frac{4 V_0(\beta_{\mathrm{liq}} - 3\alpha_g)}{\pi d_c^2}
  \end{equation}
  Maximizing linear sensitivity $S_L$ $[\si{\milli\meter\per\kelvin}]$ requires a large bulb volume $V_0$, a high-expansion working fluid $\beta_{\mathrm{liq}}$, a low-expansion borosilicate glass $\alpha_g$, and a narrow capillary bore $d_c$ (though large $V_0$ increases thermal time constant $\tau_p$).
  \item \textbf{Total vs.\ Partial Immersion \& Stem-Emergence Correction:} \idxtopic{stem-emergence correction}\idxformula{stem-emergence correction Delta_T_stem = k_g*N*(T_read - T_stem)} A \emph{partial-immersion} thermometer is calibrated with the bulb and lower stem immersed to an etched immersion ring (typically $\SI{50}{\milli\meter}$ or $\SI{76}{\milli\meter}$) while the upper stem sits in ambient air. A \emph{total-immersion} thermometer is calibrated with the entire liquid column submerged to the top of the meniscus. If a total-immersion thermometer (or a partial-immersion thermometer immersed shallower than its line, as in a $\SI{38}{\milli\meter}$-deep $\SI{100}{\milli\liter}$ cup!) has $N$ degrees (in $\si{\celsius}$ scale length) of liquid column emerging into cooler ambient stem air at $T_{\mathrm{stem}}$ while the bulb is at $T_{\mathrm{read}}$, the emergent liquid column contracts and under-reads the true bath temperature $T_{\mathrm{true}}$ by:
  \begin{equation}\label{eq:17:stem}
    \Delta T_{\mathrm{stem}} = k_g\,N\,(T_{\mathrm{read}} - T_{\mathrm{stem}}), \qquad T_{\mathrm{true}} = T_{\mathrm{read}} + k_g\,N\,(T_{\mathrm{read}} - T_{\mathrm{stem}})
  \end{equation}
  where $k_g = \beta_{\mathrm{app}} \approx \SI{0.00104}{\per\celsius}$ for organic spirit-in-glass thermometers and $k_g \approx \SI{0.00016}{\per\celsius}$ for mercury/Galinstan-in-borosilicate-glass thermometers \src{S21,S33}.
  \item \textbf{Separated Spirit Column Diagnostic \& Fix:} \idxerr{separated spirit column in thermometer} Organic spirit columns easily separate during shipping, leaving a gas bubble in the bore or a droplet clinging to the top expansion chamber (causing the thermometer to read $\SIrange{1}{5}{\celsius}$ too low!). Always inspect the entire bore before use; rejoin a separated column by cooling the bulb in an ice-salt bath so all liquid retreats into the bulb, or by gentle centrifugal swinging (never heat the bulb over an open flame, which shatters the glass expansion bulb).
\end{itemize}

\subsection{2. Negative Temperature Coefficient (NTC) Thermistors}
\idxtopic{thermistor (NTC)!Steinhart-Hart and beta model}\idxtopic{self-heating error!thermistor}
NTC thermistors are sintered metal-oxide semiconductor beads ($\mathrm{Mn_3O_4}$, $\mathrm{NiO}$, $\mathrm{Co_3O_4}$) whose charge-carrier density grows exponentially with thermal excitation across a semiconductor energy gap, causing resistance $R(T)$ to drop steeply as absolute temperature $T$ $[\si{\kelvin}]$ rises \src{S21,S33}:
\begin{equation}\label{eq:17:steinhart}
  \text{\textbf{Steinhart--Hart (3-Parameter Exact):}}\quad \frac{1}{T} = A + B\ln R + C(\ln R)^3 \quad (T\text{ in }\si{\kelvin},\ R\text{ in }\si{\ohm})
\end{equation}
\emph{(Note: The quadratic $(\ln R)^2$ term vanishes in the standard Steinhart--Hart representation, achieving $\pm\SI{0.005}{\celsius}$ fidelity across $\SIrange{0}{100}{\celsius}$).} Over a narrower $\SI{50}{\kelvin}$ window around reference state $(T_0, R_0)$ (typically $R_0 = \SI{10.00}{\kilo\ohm}$ at $T_0 = \SI{298.15}{\kelvin} = \SI{25.00}{\celsius}$), setting $C=0$ yields the \textbf{2-parameter $\beta$-equation}:
\begin{equation}\label{eq:17:beta}
  R(T) = R_0\exp\!\left[\beta\!\left(\frac{1}{T} - \frac{1}{T_0}\right)\right] \Longleftrightarrow T = \left[\frac{1}{T_0} + \frac{1}{\beta}\ln\!\left(\frac{R}{R_0}\right)\right]^{-1}, \qquad \alpha_R = \frac{1}{R}\frac{dR}{dT} = -\frac{\beta}{T^2}
\end{equation}
With material constant $\beta \in \SIrange{3400}{4100}{\kelvin}$, the fractional temperature coefficient at $T = \SI{300}{\kelvin}$ is $\alpha_R \approx -3892/300^2 = \SI{-0.0432}{\per\kelvin} = \SI{-4.32}{\percent\per\kelvin}$ --- roughly \textbf{11 times more sensitive} than a platinum RTD ($\alpha_{\mathrm{Pt}} = \SI{+0.385}{\percent\per\kelvin}$)!
\begin{itemize}
  \item \textbf{Voltage-Divider Readout \& Self-Heating Dissipation Error ($\Delta T_{\mathrm{SH}}$):} \idxformula{thermistor self-heating Delta_T_SH = I^2*R/delta_diss} In a microcontroller or digital probe readout, the thermistor $R(T)$ is placed in series with a fixed precision resistor $R_s$ across excitation voltage $V_{\mathrm{ref}}$, producing output voltage $V_{\mathrm{out}} = V_{\mathrm{ref}}\frac{R(T)}{R_s + R(T)}$ (which has maximum linearity $\frac{d^2 V_{\mathrm{out}}}{dT^2} = 0$ when $R_s = R(T_{\mathrm{mid}})\frac{\beta - 2T_{\mathrm{mid}}}{\beta + 2T_{\mathrm{mid}}} \approx 0.77\,R(T_{\mathrm{mid}})$). Because current $I = \frac{V_{\mathrm{ref}}}{R_s + R(T)}$ dissipates Joule power $P_{\mathrm{elec}} = I^2 R(T)$ inside the tiny bead, the bead warms above the surrounding fluid by:
  \begin{equation}\label{eq:17:selfheat}
    \Delta T_{\mathrm{SH}} = \frac{P_{\mathrm{elec}}}{\delta_{\mathrm{diss}}} = \frac{I^2 R(T)}{\delta_{\mathrm{diss}}} = \frac{V_{\mathrm{ref}}^2 R(T)}{\delta_{\mathrm{diss}}\left[R_s + R(T)\right]^2}
  \end{equation}
  where $\delta_{\mathrm{diss}} = h_p A_p$ $[\si{\milli\watt\per\kelvin}]$ is the probe's \textbf{dissipation constant}. Because convective transfer $h_p$ is $5\text{--}10\times$ smaller in still air than in water, $\delta_{\mathrm{diss}} \approx \SIrange{0.8}{1.5}{\milli\watt\per\kelvin}$ in still air versus $\SIrange{5.0}{10.0}{\milli\watt\per\kelvin}$ in liquid water! An improperly designed thermistor circuit carrying $I = \SI{1.0}{\milli\ampere}$ at $R = \SI{10}{\kilo\ohm}$ ($P_{\mathrm{elec}} = \SI{10}{\milli\watt}$) self-heats by $+\SI{10}{\celsius}$ in air and $+\SI{1.5}{\celsius}$ in water. Always limit thermistor excitation current to $I \le \SI{100}{\micro\ampere}$ ($P_{\mathrm{elec}} \le \SI{0.1}{\milli\watt} \implies \Delta T_{\mathrm{SH}} < \SI{0.02}{\celsius}$).
\end{itemize}

\subsection{3. Platinum Resistance Temperature Detectors (\texttt{Pt100} \& \texttt{Pt1000} RTDs)}
\idxtopic{RTD (Pt100/Pt1000)!Callendar-Van Dusen equation}\idxtopic{lead-wire resistance!3-wire and 4-wire RTD}
Pure platinum exhibits a repeatable, nearly linear positive resistance-temperature relation standardized by IEC~60751 (`Pt100`: $R_0 = \SI{100.00}{\ohm}$ at $t = \SI{0.00}{\celsius}$; `Pt1000`: $R_0 = \SI{1000.0}{\ohm}$ at $\SI{0.00}{\celsius}$) via the \textbf{Callendar--Van Dusen Equation}:
\begin{equation}\label{eq:17:cvd}
  R(t) = \begin{cases}
    R_0\left(1 + A t + B t^2\right) & \text{for }\SI{0}{\celsius} \le t \le \SI{661}{\celsius} \\[2pt]
    R_0\left[1 + A t + B t^2 + C(t - \SI{100}{\celsius})t^3\right] & \text{for }\SI{-200}{\celsius} \le t < \SI{0}{\celsius}
  \end{cases}
\end{equation}
where $t$ is temperature in $\si{\celsius}$, $A = \SI{3.9083e-3}{\per\celsius}$, $B = \SI{-5.775e-7}{\per\square\celsius}$, and $C = \SI{-4.183e-12}{\per\celsius\tothe{4}}$. Over $\SIrange{0}{100}{\celsius}$, the mean linear temperature coefficient is $\alpha_{0\text{--}100} = \frac{R(100) - R(0)}{100\,R_0} = A + 100B = \SI{0.00385055}{\per\celsius}$ ($\SI{+0.385055}{\ohm\per\celsius}$ for \texttt{Pt100}).
\begin{itemize}
  \item \textbf{IEC~60751 Accuracy Classes:} Class~AA: $\Delta t_{\mathrm{tol}} = \pm(0.10 + 0.0017|t|)\,\si{\celsius}$; Class~A: $\pm(0.15 + 0.002|t|)\,\si{\celsius}$; Class~B: $\pm(0.30 + 0.005|t|)\,\si{\celsius}$.
  \item \textbf{2-Wire Lead-Resistance Trap vs.\ 3-Wire \& 4-Wire Bridges:} \idxerr{2-wire RTD lead resistance error} Because a \texttt{Pt100} changes by only $\SI{0.385}{\ohm\per\celsius}$, two copper extension leads each of resistance $R_{\mathrm{lead}}$ in a simple 2-wire ohmmeter circuit add $2R_{\mathrm{lead}}$ to $R(t)$, creating a large positive temperature bias:
  \begin{equation}\label{eq:17:rtdlead}
    \Delta t_{\mathrm{lead,2-wire}} \approx \frac{2 R_{\mathrm{lead}}}{R_0 \alpha_{0\text{--}100}} = \frac{2 R_{\mathrm{lead}}}{\SI{0.385055}{\ohm\per\celsius}} \approx \left(\SI{5.194}{\celsius\per\ohm}\right) R_{\mathrm{lead}}
  \end{equation}
  Just $\SI{0.35}{\ohm}$ per wire produces a $+\SI{1.82}{\celsius}$ reading error! A \textbf{3-wire Wheatstone bridge} places one current-carrying lead in each adjacent bridge leg so matched lead resistances cancel to first order, while a \textbf{4-wire Kelvin measurement} uses two outer leads to supply excitation current $I$ and two separate high-impedance inner sense leads ($I_{\mathrm{sense}} \approx 0$) to measure voltage drop directly across the platinum element, eliminating lead-resistance error completely.
\end{itemize}

\subsection{4. Thermocouples (\texttt{Type K}, \texttt{T}, \texttt{J}, \texttt{E} --- NIST ITS-90 SRD~60 \src{S22})}
\idxtopic{thermocouple (Type K, T)!laws and CJC}\idxformula{Seebeck voltage Delta_V = S_AB*(T_sense - T_ref)}
When two dissimilar conductors $A$ and $B$ are joined at a sensing junction at $T_{\mathrm{sense}}$ and their open ends are held at a reference junction temperature $T_{\mathrm{ref}}$, carrier diffusion along the thermal gradient generates an open-circuit \textbf{Seebeck thermoelectric EMF} $\Delta V_{AB}$ \src{S22}:
\begin{equation}\label{eq:17:seebeck}
  \Delta V_{AB} = \int_{T_{\mathrm{ref}}}^{T_{\mathrm{sense}}} \left[S_A(T) - S_B(T)\right]dT = \int_{T_{\mathrm{ref}}}^{T_{\mathrm{sense}}} S_{AB}(T)\,dT \approx S_{AB}(\bar{T})\,(T_{\mathrm{sense}} - T_{\mathrm{ref}})
\end{equation}
\begin{itemize}
  \item \textbf{Three Fundamental Thermoelectric Laws:}
    (1)~\emph{Law of Homogeneous Circuits:} Heating a single homogeneous wire cannot produce a net thermoelectric current; voltage depends only on junction temperatures $T_{\mathrm{sense}}$ and $T_{\mathrm{ref}}$, not the temperature profile along homogeneous leads.
    (2)~\emph{Law of Intermediate Metals:} Inserting a third metal $C$ (e.g., copper voltmeter terminals or solder) into the circuit adds zero net EMF if both junctions with metal $C$ are at the same temperature.
    (3)~\emph{Law of Intermediate Temperatures (Cold-Junction Compensation, CJC):} $\Delta V(T_1, T_3) = \Delta V(T_1, T_2) + \Delta V(T_2, T_3)$. Because a digital thermocouple meter's input jack sits at ambient room temperature $T_{\mathrm{ref}} \approx \SI{22}{\celsius}$ rather than an ice bath at $\SI{0}{\celsius}$, an onboard thermistor/IC measures $T_{\mathrm{ref}}$ at the terminal block and adds the NIST ITS-90 table voltage $\Delta V(\SI{0}{\celsius}, T_{\mathrm{ref}})$ electronically \src{S22}.
  \item \textbf{Standard NIST ITS-90 Thermocouple Types at $\SI{25}{\celsius}$ \src{S22}:}
    \texttt{Type K} ($\text{Chromel}^+$ [$\mathrm{Ni\text{--}Cr}$] / $\text{Alumel}^-$ [$\mathrm{Ni\text{--}Al}$]): $S_K = \SI{40.6}{\micro\volt\per\kelvin}$ ($\SI{39.5}{\micro\volt\per\kelvin}$ at $\SI{0}{\celsius}$, $\SI{41.3}{\micro\volt\per\kelvin}$ at $\SI{50}{\celsius}$); general-purpose probe.
    \texttt{Type T} ($\text{Copper}^+$ / $\text{Constantan}^-$ [$\mathrm{Cu\text{--}Ni}$]): $S_T = \SI{40.9}{\micro\volt\per\kelvin}$; highest stability and tightest special tolerance ($\pm\SI{0.5}{\celsius}$) in wet $\SIrange{0}{100}{\celsius}$ calorimetry.
    \texttt{Type J} ($\text{Iron}^+$ / $\text{Constantan}^-$): $S_J = \SI{51.7}{\micro\volt\per\kelvin}$ (iron rusts in water; avoid in wet labs).
    \texttt{Type E} ($\text{Chromel}^+$ / $\text{Constantan}^-$): $S_E = \SI{60.9}{\micro\volt\per\kelvin}$ (highest Seebeck output and lowest thermal conductivity stem loss).
\end{itemize}

\begin{table}[h!]
\centering
\captionof{table}{Master Sensor Transduction Reference Table ($\SI{-20}{\celsius}$ to $\SI{120}{\celsius}$): \texttt{Pt100} RTD Resistance (IEC~60751), $\SI{10}{\kilo\ohm}$ NTC Thermistor ($\beta=\SI{3892}{\kelvin}$), and NIST ITS-90 Thermocouple EMF ($T_{\mathrm{ref}}=\SI{0.00}{\celsius}$) \src{S21,S22}.}\label{tab:17:sensortable}
\footnotesize
\setlength{\tabcolsep}{4.5pt}
\renewcommand{\arraystretch}{1.12}
\begin{tabularx}{\linewidth}{@{}R{0.09\linewidth} R{0.10\linewidth} R{0.14\linewidth} R{0.16\linewidth} R{0.15\linewidth} R{0.15\linewidth} >{\raggedright\arraybackslash}X@{}}
\toprule
\hdr{$t$ $[\si{\celsius}]$} & \hdr{$T$ $[\si{\kelvin}]$} & \hdr{\texttt{Pt100} $R(t)$ $[\si{\ohm}]$} & \hdr{NTC $\SI{10}{\kilo\ohm}$ $R(T)$ $[\si{\ohm}]$} & \hdr{\texttt{Type K} $E(\SI{0}{\celsius},t)$ $[\si{\milli\volt}]$} & \hdr{\texttt{Type T} $E(\SI{0}{\celsius},t)$ $[\si{\milli\volt}]$} & \hdr{Benchmark State / Calibration Note} \\
\midrule
$-20.0$ & $253.15$ & $92.16$ & $101{,}916$ & $-0.778$ & $-0.757$ & Freezer / ice-salt bath \\
\rowcolor{rowalt}
$0.0$ & $273.15$ & $100.00$ & $33{,}024$ & $0.000$ & $0.000$ & \textbf{Ice-point reference ($\delta_0$ check)} \\
$10.0$ & $283.15$ & $103.90$ & $19{,}969$ & $+0.397$ & $+0.391$ & Chilled water supply \\
\rowcolor{rowalt}
$20.0$ & $293.15$ & $107.79$ & $12{,}488$ & $+0.798$ & $+0.790$ & Standard room baseline \\
$25.0$ & $298.15$ & $109.73$ & $10{,}000$ & $+1.000$ & $+0.992$ & \textbf{NTC reference ($R_0 = \SI{10.00}{\kilo\ohm}$)} \\
\rowcolor{rowalt}
$30.0$ & $303.15$ & $111.67$ & $8{,}064$ & $+1.203$ & $+1.196$ & Mild collector warming \\
$40.0$ & $313.15$ & $115.54$ & $5{,}337$ & $+1.612$ & $+1.612$ & Typical 2027 collector $T_f$ ($h=\SI{50}{\centi\meter}$) \\
\rowcolor{rowalt}
$50.0$ & $323.15$ & $119.40$ & $3{,}613$ & $+2.023$ & $+2.036$ & High-gain collector $T_f$ ($h=\SI{40}{\centi\meter}$) \\
$60.0$ & $333.15$ & $123.24$ & $2{,}496$ & $+2.436$ & $+2.468$ & Hot water calorimetry bath \\
\rowcolor{rowalt}
$80.0$ & $353.15$ & $130.90$ & $1{,}264$ & $+3.267$ & $+3.358$ & \texttt{\#1 PET} cup $T_g$ warning zone! \\
$100.0$ & $373.15$ & $138.51$ & $682.3$ & $+4.096$ & $+4.279$ & \textbf{Normal steam point ($\SI{101.325}{\kilo\pascal}$)} \\
\rowcolor{rowalt}
$120.0$ & $393.15$ & $146.07$ & $390.0$ & $+4.920$ & $+5.228$ & Autoclave / pressurized steam \\
\bottomrule
\end{tabularx}
\end{table}

\subsection{5. Non-Contact Infrared (IR) Thermometers \& the Clear-Cup / Foil Emissivity Traps}
\idxtopic{infrared (IR) thermometer!radiometric equation}\idxerr{IR thermometer pointed through plastic cup or cling wrap}\idxerr{IR thermometer pointed at shiny aluminum foil}
A handheld IR thermometer uses a germanium or Fresnel silicon lens to focus mid-infrared thermal radiation in the atmospheric window $\lambda \in \SIrange{8}{14}{\micro\meter}$ onto a thermopile detector at ambient housing temperature $T_{\mathrm{amb}}$ \src{S12}. For an opaque gray target of true mid-IR emissivity $\varepsilon_{\mathrm{true}}$ and reflectance $\rho_{\mathrm{true}} = 1 - \varepsilon_{\mathrm{true}}$ at temperature $T_{\mathrm{obj}}$ surrounded by background walls at $T_{\mathrm{amb}}$, the total radiosity leaving the target is $J = \varepsilon_{\mathrm{true}}\sigma T_{\mathrm{obj}}^4 + (1 - \varepsilon_{\mathrm{true}})\sigma T_{\mathrm{amb}}^4$. A background-compensated IR thermometer configured with user emissivity setting $\varepsilon_{\mathrm{set}}$ (typically fixed at $\varepsilon_{\mathrm{set}} = 0.95$) solves $\varepsilon_{\mathrm{set}}\sigma T_{\mathrm{read}}^4 + (1 - \varepsilon_{\mathrm{set}})\sigma T_{\mathrm{amb}}^4 = J$, yielding the exact apparent temperature read by the IR gun:
\begin{equation}\label{eq:17:ir}
  \boxed{T_{\mathrm{read}} = \left[T_{\mathrm{amb}}^4 + \frac{\varepsilon_{\mathrm{true}}}{\varepsilon_{\mathrm{set}}}\left(T_{\mathrm{obj}}^4 - T_{\mathrm{amb}}^4\right)\right]^{1/4} \quad (\text{all }T\text{ in }\si{\kelvin})}
\end{equation}

\begin{trapbox}[Four Critical Infrared Thermometer Traps in 2027 Thermodynamics Testing]
\begin{enumerate}
  \item \textbf{Trap 1 --- Shooting Through the Clear \texttt{\#1 PET} / \texttt{\#5 PP} Plastic Cup Sidewall or Cling-Wrap Glazing:} \idxerr{shooting IR gun through plastic wrap or cup wall}
  Because water and plastics are visually transparent at $\lambda = \SIrange{0.4}{0.7}{\micro\meter}$, students falsely assume an IR thermometer can ``see through'' a clear $\SI{9}{oz}$ plastic cup wall or a plastic cling-wrap cover to read the water inside. In the $\SIrange{8}{14}{\micro\meter}$ thermal-IR band, glass and clear plastics ($\ge\SI{0.3}{\milli\meter}$ PET/PP) are \textbf{virtually opaque} ($\tau_{\mathrm{MIR}} < 0.05\text{--}0.25$, $\varepsilon_{\mathrm{MIR}} \approx 0.90\text{--}0.94$) \src{S12}! Shooting an IR gun at the cup sidewall or through an unremoved top plastic film measures the \textbf{cooler exterior plastic temperature}, not the water! Under FAQ 10/21/2026 (items on the cup) and the ES Guide \src{S02,S05}, any team-placed film or lid \textbf{must be completely removed from the cup} during the 1-minute Final Measurement window so the ES thermometer views or contacts the bare top water surface directly.
  \item \textbf{Trap 2 --- Shooting at Bare Liquid Water Surface (Why Top-Down IR Works Well on Open Water):}
  In the $\SIrange{8}{14}{\micro\meter}$ band, liquid water has an absorption coefficient $\mu_a > \SI{600}{\per\centi\meter}$ (penetration depth $\delta_p < \SI{17}{\micro\meter}$) and a hemispherical emissivity of $\varepsilon_w = \mathbf{0.96}$ at $\SI{40}{\celsius}$ (normal-incidence $\varepsilon_{w,n} = 0.985$, Lienhard Table~10.1 p.~546 \src{S12}). Thus, an IR gun preset to $\varepsilon_{\mathrm{set}} = 0.95$ aimed vertically downward at bare water reads within $+\SI{0.2}{\celsius}$ of the true surface skin temperature (which is slightly cooled by evaporation once uncovered).
  \item \textbf{Trap 3 --- Shooting at Polished Metals or Aluminum Foil ($\varepsilon_{\mathrm{true}} = 0.04$):}
  On shiny aluminum foil ($\varepsilon_{\mathrm{true}} = 0.04$, $\rho_{\mathrm{true}} = 0.96$), $96\%$ of the IR radiation entering the gun is reflected room background! By \eref{eq:17:ir}, a $\SI{44.0}{\celsius}$ aluminum cone liner in a $\SI{22.0}{\celsius}$ room reads only $\SI{23.0}{\celsius}$ on an IR gun with $\varepsilon_{\mathrm{set}} = 0.95$ --- a $-\SI{21.0}{\celsius}$ error! To measure a shiny metal surface with an IR gun, apply a patch of matte black electrical tape or flat black paint ($\varepsilon = 0.95$) and target the patch.
  \item \textbf{Trap 4 --- Distance-to-Spot Ratio ($D:S$) Field-of-View Spillover:} \idxtopic{distance-to-spot ratio (D:S)}
  An IR gun with optical ratio $D:S = 12:1$ held at distance $d$ views a circular spot of diameter $D_{\mathrm{spot}} = d / 12$. Because the $\SI{9}{oz}$ cup water surface has diameter $D_{\mathrm{water}} \approx \SI{6.8}{\centi\meter}$, holding a $12:1$ IR gun farther than $d_{\max} = 12 \times (\SI{4.5}{\centi\meter}\text{ safe central zone}) = \SI{54}{\centi\meter}$ causes the sensor spot to spill over onto the cooler cup rim and table, averaging down $T_f$! Always hold the IR gun at $d = \SIrange{15}{30}{\centi\meter}$ ($D_{\mathrm{spot}} = \SIrange{1.3}{2.5}{\centi\meter}$) centered directly over the cup.
\end{enumerate}
\end{trapbox}

\begin{workedexample}[Liquid-in-Glass Capillary Design \& Partial-Immersion Stem Correction]{Rule 4.f \& 5.c.i}{ex:17.1}
A spirit-in-glass laboratory thermometer uses ethanol ($\beta_{\mathrm{liq}} = \SI{1.120e-3}{\per\kelvin}$) in a borosilicate glass bulb ($\alpha_g = \SI{3.30e-6}{\per\kelvin}$) of internal volume $V_0 = \SI{250.0}{\cubic\milli\meter}$ connected to a capillary bore of diameter $d_c = \SI{0.400}{\milli\meter}$. (a)~Calculate the linear scale sensitivity $S_L = \Delta L/\Delta T$ in $\si{\milli\meter\per\kelvin}$. (b)~During a hot-water bath check, the thermometer reads $T_{\mathrm{read}} = \SI{65.00}{\celsius}$, but an emergent column of length equivalent to $N = \SI{45.0}{\celsius}$ protrudes above the immersion line into ambient air at $T_{\mathrm{stem}} = \SI{23.0}{\celsius}$ ($k_g = \SI{0.00104}{\per\celsius}$). Find the true bath temperature $T_{\mathrm{true}}$.
\par\smallskip\noindent
\textbf{Solution.}
(a)~Apparent expansion coefficient: $\beta_{\mathrm{app}} = \beta_{\mathrm{liq}} - 3\alpha_g = \num{1.120e-3} - 3(\num{3.30e-6}) = \SI{1.1101e-3}{\per\kelvin}$. Capillary cross-sectional area: $A_c = \frac{\pi}{4}(\SI{0.400}{\milli\meter})^2 = \SI{0.125664}{\square\milli\meter}$. By \eref{eq:17:lig}:
$$S_L = \frac{\Delta L}{\Delta T} = \frac{V_0 \beta_{\mathrm{app}}}{A_c} = \frac{(\SI{250.0}{\cubic\milli\meter})(\SI{1.1101e-3}{\per\kelvin})}{\SI{0.125664}{\square\milli\meter}} = \mathbf{\SI{2.2085}{\milli\meter\per\kelvin}}$$
(b)~By \eref{eq:17:stem}, the emergent column correction is $\Delta T_{\mathrm{stem}} = k_g N (T_{\mathrm{read}} - T_{\mathrm{stem}}) = (\num{0.00104})(45.0)(65.00 - 23.0) = +\SI{1.9656}{\celsius}$, so:
$$T_{\mathrm{true}} = 65.00 + 1.9656 = \mathbf{\SI{66.97}{\celsius}}$$
\verifiedby{Independent volumetric balance $\Delta V_{\mathrm{app}}(\SI{1}{\kelvin}) = \SI{0.277525}{\cubic\milli\meter} = A_c(\SI{2.20847}{\milli\meter})$ and Python check \texttt{sec17:ex17.1} ($\Delta T_{\mathrm{stem}} = +\SI{1.9656}{\celsius}$).}
\end{workedexample}

\begin{workedexample}[NTC Thermistor $\beta$-Model Inversion, Sensitivity \& Self-Heating Audit]{Rule 4.f}{ex:17.2}
An Event Supervisor's digital probe uses a $\SI{10.00}{\kilo\ohm}$ NTC thermistor ($R_0 = \SI{10000}{\ohm}$ at $T_0 = \SI{25.00}{\celsius} = \SI{298.15}{\kelvin}$, $\beta = \SI{3892}{\kelvin}$, dissipation constant in water $\delta_{\mathrm{diss}} = \SI{6.50}{\milli\watt\per\kelvin}$) connected in a voltage divider with $R_s = \SI{10.00}{\kilo\ohm}$ and $V_{\mathrm{ref}} = \SI{3.30}{\volt}$. When immersed in the heated cup, the thermistor resistance is $R = \SI{4912.0}{\ohm}$. Find: (a)~the water temperature in $\si{\celsius}$, (b)~the fractional sensitivity $\alpha_R$ and $\frac{dR}{dT}$ at that temperature, and (c)~the Joule self-heating bias $\Delta T_{\mathrm{SH}}$.
\par\smallskip\noindent
\textbf{Solution.}
(a)~From \eref{eq:17:beta}:
$$\frac{1}{T} = \frac{1}{\SI{298.15}{\kelvin}} + \frac{1}{\SI{3892}{\kelvin}}\ln\!\left(\frac{4912.0}{10000.0}\right) = \num{3.3540164e-3} - \num{1.826577e-4} = \SI{3.1713587e-3}{\per\kelvin} \implies T = \SI{315.3223}{\kelvin} = \mathbf{\SI{42.17}{\celsius}}$$
(b)~Fractional sensitivity: $\alpha_R = -\frac{\beta}{T^2} = -\frac{3892}{(315.3223)^2} = \mathbf{\SI{-0.03914}{\per\kelvin}} = \mathbf{\SI{-3.914}{\percent\per\kelvin}}$, so $\frac{dR}{dT} = \alpha_R R = (\num{-0.039144})(4912.0) = \mathbf{\SI{-192.3}{\ohm\per\kelvin}}$.
(c)~Divider current: $I = \frac{3.30}{10000 + 4912.0} = \SI{0.22130}{\milli\ampere}$. Electrical dissipation: $P_{\mathrm{elec}} = I^2 R = (\num{2.2130e-4})^2(4912.0) = \SI{0.24056}{\milli\watt}$. By \eref{eq:17:selfheat}, $\Delta T_{\mathrm{SH}} = \frac{\SI{0.24056}{\milli\watt}}{\SI{6.50}{\milli\watt\per\kelvin}} = \mathbf{+\SI{0.0370}{\celsius}}$ (well below the $\SI{0.1}{\celsius}$ display resolution; however, in still air where $\delta_{\mathrm{diss}} \approx \SI{1.0}{\milli\watt\per\kelvin}$, $\Delta T_{\mathrm{SH}}$ rises to $+\SI{0.24}{\celsius}$!).
\verifiedby{Forward substitution $R(\SI{315.3223}{\kelvin}) = 10000\exp[3892(1/315.3223 - 1/298.15)] = \SI{4912.00}{\ohm}$ and Python check \texttt{sec17:ex17.2}.}
\end{workedexample}

\begin{workedexample}[\texttt{Pt100} RTD Lead-Resistance Bias \& IR Pyrometer Emissivity Audit]{Rule 4.f}{ex:17.3}
(a)~A 2-wire \texttt{Pt100} RTD ($R_0 = \SI{100.00}{\ohm}$, $A = \SI{3.9083e-3}{\per\celsius}$, $B = \SI{-5.775e-7}{\per\square\celsius}$) has two extension leads of resistance $R_{\mathrm{lead}} = \SI{0.350}{\ohm}$ each. If the meter reads total loop resistance $R_{\mathrm{meas}} = \SI{117.47}{\ohm}$, find the true water temperature $t_{\mathrm{true}}$ vs.\ the uncorrected 2-wire reading $t_{\mathrm{app}}$.
(b)~An IR thermometer ($\varepsilon_{\mathrm{set}} = 0.95$, ambient room $T_{\mathrm{amb}} = \SI{22.00}{\celsius} = \SI{295.15}{\kelvin}$) is aimed at bare water ($\varepsilon_w = 0.96$) at $T_{\mathrm{obj}} = \SI{44.00}{\celsius} = \SI{317.15}{\kelvin}$, and then at the collector's shiny aluminum-foil cone ($\varepsilon_{\mathrm{Al}} = 0.04$) also at $\SI{44.00}{\celsius}$. Compute both IR readings.
\par\smallskip\noindent
\textbf{Solution.}
(a)~Subtracting the two-lead loop resistance $2R_{\mathrm{lead}} = \SI{0.700}{\ohm}$ gives true RTD resistance $R_{\mathrm{true}} = 117.47 - 0.700 = \SI{116.77}{\ohm}$. Solving the Callendar--Van Dusen quadratic \eref{eq:17:cvd}, $100.00(1 + \num{3.9083e-3}t - \num{5.775e-7}t^2) = R$:
$$t = \frac{-A + \sqrt{A^2 - 4B(1 - R/R_0)}}{2B} \implies t_{\mathrm{true}}(\SI{116.77}{\ohm}) = \mathbf{\SI{43.18}{\celsius}}, \qquad t_{\mathrm{app}}(\SI{117.47}{\ohm}) = \mathbf{\SI{45.00}{\celsius}}$$
The 2-wire lead bias is $t_{\mathrm{app}} - t_{\mathrm{true}} = \mathbf{+\SI{1.815}{\celsius}}$ (eliminated by a 3-wire or 4-wire Kelvin probe).
\par\noindent
(b)~Using the background-compensated IR radiometer equation \eref{eq:17:ir} with $T_{\mathrm{amb}} = \SI{295.15}{\kelvin}$ and $T_{\mathrm{obj}} = \SI{317.15}{\kelvin}$:
$$T_{\mathrm{read,water}} = \left[295.15^4 + \frac{0.96}{0.95}\left(317.15^4 - 295.15^4\right)\right]^{1/4} - 273.15 = 317.358 - 273.15 = \mathbf{\SI{44.21}{\celsius}} \quad (+\SI{0.21}{\celsius}\text{ offset})$$
$$T_{\mathrm{read,foil}} = \left[295.15^4 + \frac{0.04}{0.95}\left(317.15^4 - 295.15^4\right)\right]^{1/4} - 273.15 = 296.180 - 273.15 = \mathbf{\SI{23.03}{\celsius}} \quad (-\SI{20.97}{\celsius}\text{ error!})$$
\verifiedby{Quadratic verification $100(1 + \num{3.9083e-3}(43.1842) - \num{5.775e-7}(43.1842)^2) = \SI{116.770}{\ohm}$ and Python check \texttt{sec17:ex17.3}.}
\end{workedexample}

% ============================================================================
\section{Metrology Fundamentals, Parallax, Balances, Timers \& Volumetric Repeatability}\label{sec:17.2}
\idxtopic{accuracy vs precision vs resolution}\idxtopic{parallax error}\idxtopic{volumetric dispensing (100 mL)}

\subsection{1. Resolution vs.\ Precision vs.\ Accuracy vs.\ Sensitivity}
\begin{itemize}
  \item \textbf{Resolution ($\Delta_{\mathrm{res}}$):} The smallest increment distinguishable on a display or scale (e.g., $\Delta_{\mathrm{res}} = \SI{0.1}{\celsius}$ required by Rule~\texttt{4.f}). On an analog scale with graduation spacing $\Delta_{\mathrm{grad}}$, human interpolation resolves $\Delta_{\mathrm{res}} \approx \Delta_{\mathrm{grad}}/5$ to $\Delta_{\mathrm{grad}}/2$, whereas a digital display has hard quantization $\Delta_{\mathrm{res}} = 1\text{ LSD}$ (least significant digit), creating a rectangular Type~B rounding uncertainty $u_{B,\mathrm{res}} = \frac{\Delta_{\mathrm{res}}}{2\sqrt{3}} \approx 0.2887\,\Delta_{\mathrm{res}}$ (\sref{sec:17.6}) \src{S21}.
  \item \textbf{Precision (Repeatability / Random Scatter):} Closeness of agreement among independent repeated measurements under identical conditions, quantified by sample standard deviation $s$ and Type~A standard uncertainty of the mean $u_A = s/\sqrt{n}$. High precision means small random noise, even if all readings share a constant offset.
  \item \textbf{Accuracy (Trueness / Low Systematic Bias):} Closeness of agreement between the mean measured value and the true NIST/SI reference value. Systematic bias $\delta_{\mathrm{sys}} = \bar{T}_{\mathrm{meas}} - T_{\mathrm{true}}$ is reduced only by physical calibration (\sref{sec:17.5}) or differential cancellation.
  \item \textbf{Why Rule~\texttt{4.f} (``Same Thermometer for Start \& Final Reading'') Cancels Zero-Offset Bias in $\mathrm{TG}$:} \idxtopic{Rule 4.f!systematic bias cancellation}
  Suppose the ES thermometer has a constant additive zero-offset bias $\delta_{\mathrm{sys}}$ (e.g., reads $+\SI{0.6}{\celsius}$ high across $\SIrange{20}{50}{\celsius}$): $T_{\mathrm{start,read}} = T_{\mathrm{start,true}} + \delta_{\mathrm{sys}}$ and $T_{f,\mathrm{read}} = T_{f,\mathrm{true}} + \delta_{\mathrm{sys}}$.
  \begin{enumerate}
    \item In the \textbf{Temperature Score ($\mathrm{TS}$, 30 pts, Rule~\texttt{6.e})}, which depends on Temperature Gain $\mathrm{TG} = T_{f,\mathrm{read}} - T_{\mathrm{start,read}} = (T_{f,\mathrm{true}} + \delta_{\mathrm{sys}}) - (T_{\mathrm{start,true}} + \delta_{\mathrm{sys}}) = \mathrm{TG}_{\mathrm{true}}$, the additive zero offset $\delta_{\mathrm{sys}}$ \textbf{cancels identically}!
    \item In the \textbf{Prediction Score ($\mathrm{PS}$, 20 pts, Rule~\texttt{6.f})}, $\mathrm{PE} = 1 - \frac{|T_{f,\mathrm{read}} - T_{\mathrm{pred}}|}{T_{f,\mathrm{read}}}$. Because the ES announces $T_{\mathrm{start,read}}$ measured with that \textbf{same} offset thermometer (\texttt{Rule 5.II.c}) and your team computes $T_{\mathrm{pred}} = T_{\mathrm{start,read}} + \widehat{\mathrm{TG}}$, the numerator $|T_{f,\mathrm{read}} - T_{\mathrm{pred}}| = |(T_{f,\mathrm{true}} + \delta_{\mathrm{sys}}) - (T_{\mathrm{start,true}} + \delta_{\mathrm{sys}} + \widehat{\mathrm{TG}})| = |\mathrm{TG}_{\mathrm{true}} - \widehat{\mathrm{TG}}|$ \textbf{also cancels $\delta_{\mathrm{sys}}$ identically}! Always anchor your prediction equation directly to the ES-announced $T_{\mathrm{start}}$ (\sref{sec:16}).
  \end{enumerate}
\end{itemize}

\subsection{2. Meniscus Parallax, Digital Balances, Timers \& $\SI{100}{\milli\liter}$ Water Dispensing Repeatability}
\idxtopic{meniscus reading!concave vs convex}\idxtopic{buoyancy correction!gravimetric weighing}
\begin{table}[h!]
\centering
\captionof{table}{Laboratory Metrology Standards for Volume, Mass, Time, and Meniscus Reading in Calibration Labs \src{S05,S09,S21}.}\label{tab:17:labgear}
\footnotesize
\setlength{\tabcolsep}{3.5pt}
\renewcommand{\arraystretch}{1.15}
\begin{tabularx}{\linewidth}{@{}>{\raggedright\arraybackslash}p{0.21\linewidth} >{\raggedright\arraybackslash}p{0.34\linewidth} C{0.14\linewidth} >{\raggedright\arraybackslash}X@{}}
\toprule
\hdr{Measurement / Instrument} & \hdr{Governing Physics \& Systematic Bias Formula} & \hdr{Typical Uncertainty} & \hdr{Best-Practice Protocol for 2027 Device Calibration} \\
\midrule
\textbf{Meniscus \& Parallax} \newline (Graduated cylinder \& spirit thermometer) &
Line-of-sight angle $\theta$ off horizontal across bore-to-scale gap $d_g$ shifts reading by $\Delta y_{\mathrm{par}} = d_g \tan\theta \approx d_g(\Delta z_{\mathrm{eye}} / L_{\mathrm{eye}})$. &
$\pm\SIrange{0.2}{0.5}{\milli\liter}$ or $\pm\SI{0.3}{\celsius}$ if viewed at angle &
Water \& ethanol wet glass (adhesion $>$ cohesion $\implies$ \textbf{concave meniscus}: read \textbf{bottom} of trough at eye level). Mercury does not wet glass (\textbf{convex}: read top of dome). \\
\rowcolor{rowalt}
\textbf{Volumetric $\SI{100}{\milli\liter}$ Dispensing} (\texttt{Rule 3.a, 5.II.c}) &
Thermal capacity of water $C_w = \rho_w(T)\,V_w\,c_{p,w}$ scales linearly with volume: $\dfrac{\delta(\mathrm{TG})}{\mathrm{TG}} \approx -\dfrac{C_w}{C_{\mathrm{eff}}}\dfrac{\delta V_w}{V_w}$ ($\approx -\SI{0.20}{\celsius}$ per $+\SI{1.0}{\milli\liter}$ error at $\mathrm{TG}=\SI{20}{\celsius}$!). &
Beaker: $\pm\SI{5}{\milli\liter}$ \newline Class B cyl: $\pm\SI{1.0}{\milli\liter}$ \newline Class A cyl: $\pm\SI{0.5}{\milli\liter}$ \newline Gravimetric: $\pm\SI{0.05}{\milli\liter}$ &
During pre-tournament DOE calibration (\sref{sec:16}), always dispense the $\SI{100.0}{\milli\liter}$ water sample \textbf{gravimetrically on a $\SI{0.01}{\gram}$ digital balance} so volume scatter never corrupts your calibration curve! \\
\textbf{Gravimetric Mass-to-Volume \& Air Buoyancy} \src{S09,S21} &
True mass $m_{\mathrm{true}} = m_{\mathrm{bal}}\dfrac{1 - \rho_{\mathrm{air}}/\rho_{\mathrm{cal}}}{1 - \rho_{\mathrm{air}}/\rho_w} \approx m_{\mathrm{bal}}\left[1 + \rho_{\mathrm{air}}\!\left(\dfrac{1}{\rho_w} - \dfrac{1}{\rho_{\mathrm{cal}}}\right)\right] = 1.00106\,m_{\mathrm{bal}}$ ($\rho_{\mathrm{air}} = \SI{1.20}{\kilogram\per\cubic\meter}$, steel calibration weight $\rho_{\mathrm{cal}} = \SI{8000}{\kilogram\per\cubic\meter}$). &
Balance: $\pm\SI{0.01}{\gram}$ \newline Buoyancy + $\rho_w(T)$ shift: $\SI{0.33}{\gram}$ at $\SI{22}{\celsius}$ &
At $T = \SI{22.0}{\celsius}$, $\rho_w = \SI{0.99777}{\gram\per\milli\liter}$ \src{S09}; accounting for air buoyancy ($1.00106$), $\SI{100.00}{\milli\liter}$ of water weighs exactly $m_{\mathrm{bal}} = \dfrac{\SI{99.777}{\gram}}{1.00106} = \mathbf{\SI{99.67}{\gram}}$ on a tared digital scale (at $\SI{20.0}{\celsius}$, $\SI{99.72}{\gram}$; at $\SI{25.0}{\celsius}$, $\SI{99.60}{\gram}$). \\
\rowcolor{rowalt}
\textbf{Digital Stopwatch \& AC Line Voltage} &
Heating gain scales as $\dfrac{\partial T_f}{\partial t} \approx \SI{1.5}{\celsius\per\minute} = \SI{0.025}{\celsius\per\second}$; incandescent lamp power scales with AC mains voltage as $P_{\mathrm{bulb}} \propto (V_{\mathrm{AC}}/V_0)^{1.54}$. &
Human reaction: $\pm\SI{0.2}{\second}$ ($\pm\SI{0.005}{\celsius}$) \newline Mains voltage: $\pm\SI{3}{\percent}$ &
Stopwatch reaction error is negligible, but a $\SI{3}{\percent}$ building AC line voltage sag ($\SI{120}{\volt} \to \SI{116.4}{\volt}$) reduces heat-lamp output by $1.54 \times \SI{3}{\percent} = \SI{4.6}{\percent}$. Log Kill-A-Watt power meter readings during calibration runs. \\
\bottomrule
\end{tabularx}
\end{table}

% ============================================================================
\section{Dynamic Probe Response Time ($\tau_p$), First-Order Thermal Lag \& the $5\tau_p$ Settling Rule}\label{sec:17.3}
\idxtopic{probe response time (tau_p)}\idxtopic{5-tau settling rule}\idxformula{thermometer step response T_p(t) = T_w - (T_w - T_p0)*exp(-t/tau_p)}

When a thermometer probe of sensing-tip mass $m_p$, specific heat $c_p$, wetted surface area $A_p$, and internal thermal conductivity $k_p$ (with probe Biot number $\mathrm{Bi}_p = \frac{h_p r_p}{2 k_p} < 0.1$) is transferred at $t = 0$ from ambient room air at $T_{p,0}$ into liquid water at temperature $T_w$, convective heat transfer across the probe boundary layer ($h_p$) governs the probe reading $T_p(t)$ via a first-order linear energy balance \src{S12}:
\begin{equation}\label{eq:17:probeode}
  m_p c_p \frac{dT_p}{dt} = h_p A_p \left[T_w - T_p(t)\right] \implies \tau_p \frac{dT_p}{dt} + T_p(t) = T_w, \qquad \tau_p = \frac{m_p c_p}{h_p A_p} = R_{\mathrm{th},p}\,C_p
\end{equation}

\begin{formulacard}[First-Order Probe Step Response, $5\tau_p$ Settling Rule, Ramp Lag \& Probe Thermal Drain]
\begin{enumerate}
  \item \textbf{Step-Immersion Response \& Dynamic Fractional Lag Error:}
  \begin{equation}\label{eq:17:probestep}
    T_p(t) = T_w - \left(T_w - T_{p,0}\right)e^{-t/\tau_p}, \qquad E_{\mathrm{frac}}(t) = \frac{T_w - T_p(t)}{T_w - T_{p,0}} = e^{-t/\tau_p}, \qquad f_{\mathrm{settled}}(t) = 1 - e^{-t/\tau_p}
  \end{equation}
  \begin{center}
  \footnotesize
  \setlength{\tabcolsep}{5pt}
  \begin{tabular}{@{}lcccccc@{}}
  \toprule
  \textbf{Elapsed Time $t/\tau_p$} & $1\tau_p$ & $2\tau_p$ & $3\tau_p$ & $4\tau_p$ & $\mathbf{5\tau_p}$ \textbf{(Rule)} & $6\tau_p$ \\
  \midrule
  \textbf{Fraction Settled $1 - e^{-t/\tau_p}$} & $\SI{63.21}{\percent}$ & $\SI{86.47}{\percent}$ & $\SI{95.02}{\percent}$ & $\SI{98.17}{\percent}$ & $\mathbf{\SI{99.326}{\percent}}$ & $\SI{99.752}{\percent}$ \\
  \textbf{Remaining Error $e^{-t/\tau_p}$} & $\SI{36.79}{\percent}$ & $\SI{13.53}{\percent}$ & $\SI{4.98}{\percent}$ & $\SI{1.83}{\percent}$ & $\mathbf{\SI{0.674}{\percent}}$ & $\SI{0.248}{\percent}$ \\
  \textbf{Lag Error for $\Delta T = \SI{25.0}{\celsius}$ Step} & $-\SI{9.20}{\celsius}$ & $-\SI{3.38}{\celsius}$ & $-\SI{1.24}{\celsius}$ & $-\SI{0.46}{\celsius}$ & $\mathbf{-\SI{0.17}{\celsius}}$ & $-\SI{0.06}{\celsius}$ \\
  \bottomrule
  \end{tabular}
  \end{center}
  \item \textbf{Two-Point Algebraic Extraction of Probe Time Constant $\tau_p$ (Zero Calculus Required):} \idxformula{probe time constant from two readings tau_p = (t2 - t1)/ln((T_w - T_p1)/(T_w - T_p2))}
  From any two probe readings $T_p(t_1)$ and $T_p(t_2)$ during step immersion into a constant-temperature bath $T_w$:
  \begin{equation}\label{eq:17:taup2pt}
    \tau_p = \frac{t_2 - t_1}{\ln\!\left(\dfrac{T_w - T_p(t_1)}{T_w - T_p(t_2)}\right)}
  \end{equation}
  \item \textbf{Dynamic Ramp Lag During Continuous Heating ($T_w(t) = T_{w,0} + r\,t$):} \idxtopic{dynamic ramp lag of thermometer}\idxformula{ramp lag Delta_T_lag = -r*tau_p}
  If a thermometer probe remains immersed while water warms at rate $r = \frac{dT_w}{dt}$ $[\si{\celsius\per\second}]$, the exact steady-ramp solution to $\tau_p \frac{dT_p}{dt} + T_p = T_{w,0} + rt$ is:
  \begin{equation}\label{eq:17:ramplag}
    T_p(t) = T_w(t) - r\tau_p\left(1 - e^{-t/\tau_p}\right) \xrightarrow{t \ge 5\tau_p} T_p(t) - T_w(t) = -r\,\tau_p
  \end{equation}
  \emph{(During continuous heating at $r = \SI{2.4}{\celsius\per\minute} = \SI{0.040}{\celsius\per\second}$, a slow spirit thermometer with $\tau_p = \SI{10}{\second}$ constantly trails the true water temperature by $-r\tau_p = -\SI{0.40}{\celsius}$!)}
  \item \textbf{Calorimetric Probe Heat-Capacity Perturbation ($\Delta T_{\mathrm{drain}}$):} \idxtopic{probe heat-capacity drain}\idxformula{probe thermal drain Delta_T_drain = -C_p*(T_w0 - T_p0)/(C_w + C_p)}
  When a cold probe of effective submerged heat capacity $C_p = m_p c_p$ initially at room temperature $T_{p,0}$ is dipped into an unstirred top water layer of effective heat capacity $C_{\mathrm{top}}$ initially at $T_{w,0}$, thermal equilibration cools the local water layer to:
  \begin{equation}\label{eq:17:probedrain}
    T_{\mathrm{eq}} = \frac{C_{\mathrm{top}} T_{w,0} + C_p T_{p,0}}{C_{\mathrm{top}} + C_p} \implies \Delta T_{\mathrm{drain}} = T_{\mathrm{eq}} - T_{w,0} = -\frac{C_p}{C_{\mathrm{top}} + C_p}\left(T_{w,0} - T_{p,0}\right)
  \end{equation}
\end{enumerate}
\end{formulacard}

\begin{center}
\begin{tikzpicture}
\begin{axis}[
    width=0.92\linewidth,
    height=5.2cm,
    font=\scriptsize,
    xlabel={Dimensionless Immersion Time $t / \tau_p$ (for $\tau_p = \SI{4.0}{\second}$, $1\tau_p = \SI{4}{\second}$ and $5\tau_p = \SI{20}{\second}$)},
    ylabel={Normalized Response / Fractional Error},
    xmin=0, xmax=6.0,
    ymin=0, ymax=1.08,
    xtick={0,1,2,3,4,5,6},
    xticklabels={$0$, $1\tau_p\ (\SI{4}{\second})$, $2\tau_p\ (\SI{8}{\second})$, $3\tau_p\ (\SI{12}{\second})$, $4\tau_p\ (\SI{16}{\second})$, $\mathbf{5\tau_p\ (\SI{20}{\second})}$, $6\tau_p\ (\SI{24}{\second})$},
    ytick={0,0.2,0.4,0.632,0.8,0.95,1.0},
    yticklabels={$0\%$, $20\%$, $40\%$, $63.2\%$, $80\%$, $95.0\%$, $100\%$},
    grid=both,
    grid style={gray!25, dashed},
    legend style={at={(0.52,0.52)}, anchor=west, font=\scriptsize, fill=white, draw=navy},
     samples=36,
    domain=0:6
]
  % Settling curve 1 - exp(-x)
  \addplot[navy, very thick] {1 - exp(-x)};
  \addlegendentry{Probe Step Response $\dfrac{T_p(t) - T_{p,0}}{T_{\mathrm{true}} - T_{p,0}} = 1 - e^{-t/\tau_p}$}
  % Error curve exp(-x)
  \addplot[rust, very thick, dashed] {exp(-x)};
  \addlegendentry{Fractional Lag Error $E_{\mathrm{frac}}(t) = \dfrac{T_{\mathrm{true}} - T_p(t)}{\Delta T} = e^{-t/\tau_p}$}
  % Initial tangent line y = x
  \addplot[teal, semithick, dotted, domain=0:1.05] {x};
  \addlegendentry{Initial Slope Tangent at $t=0$: $(d\theta/dt)_0 = 1/\tau_p$}

  % Key markers at 1 tau, 3 tau, 5 tau
  \addplot[only marks, mark=*, mark size=2pt, navy] coordinates {(1, 0.6321) (2, 0.8647) (3, 0.9502) (5, 0.9933)};
  \addplot[only marks, mark=square*, mark size=1.8pt, rust] coordinates {(1, 0.3679) (2, 0.1353) (3, 0.0498) (5, 0.0067)};

  \draw[navy, densely dashed] (axis cs:1,0) -- (axis cs:1,0.6321);
  \draw[navy, densely dashed] (axis cs:3,0) -- (axis cs:3,0.9502);
  \draw[rust, thick] (axis cs:5,0) -- (axis cs:5,1.02);

  \node[anchor=south west, font=\scriptsize, fill=white, inner sep=1pt] at (axis cs:1.05,0.60) {$1\tau_p$: $63.2\%$ ($36.8\%$ err)};
  \node[anchor=south east, font=\scriptsize, fill=white, inner sep=1pt] at (axis cs:3.0,0.96) {$3\tau_p$: $95.0\%$ ($5.0\%$ err)};
  \node[anchor=north east, font=\scriptsize\bfseries, text=rust, fill=boxbg, draw=rust, inner sep=2pt] at (axis cs:4.92,0.91) {$5\tau_p$ Settling Threshold: $99.33\%$ settled ($0.67\%$ err $< \SI{0.17}{\celsius}$)};
\end{axis}
\end{tikzpicture}
\captionof{figure}{First-order thermometer probe step response $1 - e^{-t/\tau_p}$ and fractional dynamic lag error $e^{-t/\tau_p}$ versus dimensionless immersion time $t/\tau_p$. A probe with $\tau_p = \SI{4.0}{\second}$ must remain immersed for $\ge 5\tau_p = \SI{20.0}{\second}$ to reduce dynamic lag below $0.67\%$ of $\Delta T$ \src{S12,S21}.}\label{fig:17:thermometer-response}
\idxfig{First-Order Thermometer Probe Step Response and 5-tau Settling Rule}
\end{center}

% ============================================================================
\section{Immersion Depth, Stem Conduction \& Radial Stratification Bias}\label{sec:17.4}
\idxtopic{stratification bias!probe depth and radial placement}\idxdev{top-layer probe placement (FAQ 10/06/2026)}\idxtopic{FAQs on stirring and probe depth!probe placement}

\begin{rulebox}[FAQs 09/30, 10/01 and 10/06/2026: Unstirred Cup \& Top-Layer Measurement \src{S02,S05}]
\begin{itemize}
  \item \textbf{FAQ 10/06/2026 (also 09/30 and 10/01/2026):} \emph{``No, the water must not be swirled, shaken or stirred before taking the temperature.''}
  \item \textbf{FAQ 10/06/2026 (probe depth):} \emph{``This will be up to the local Event Supervisor to determine the temperature measuring protocol \ldots\ It is recommended that the ES should measure the temperature of the top layer of water near the center of the cup. If the thermometer has a probe, the ES should insert it far enough into the water to take an accurate measurement, but no farther.''}
\end{itemize}
\end{rulebox}

Why does probe placement inside the $\SI{100}{\milli\liter}$ water column ($\SI{38}{\milli\meter}$ total fill depth in a standard $\SI{9}{oz}$ clear cup) shift the measured temperature $T_f$ by up to about $\SI{19}{\celsius}$ (\fref{fig:17:probe-depth-bias})?
\begin{enumerate}
  \item \textbf{Vertical Optical Deposition \& Hydrostatic Stability ($z$-Axis Bias):}
  As established by the Hale \& Querry (1973) and Kou et al.\ (1993) liquid-water near-IR absorption spectra (\sref{sec:14}, \sref{sec:15.3}) \src{S19,S20}, $>66\%$ of the $\SI{2400}{\kelvin}$ tungsten heat lamp's transmitted radiant energy is absorbed directly inside the \textbf{top $z = \SI{5}{\milli\meter}$} of water ($>72\%$ in the top $\SI{10}{\milli\meter}$). Because liquid water above $\SI{4}{\celsius}$ expands when heated ($\partial\rho/\partial T < 0$), the hot top water layer is less dense than the cooler water beneath it ($\partial\rho/\partial z > 0$ downward). This gravitationally stable density profile completely suppresses vertical natural-convection overturning ($\mathrm{Ra} < 0$), trapping heat in the upper $\SIrange{5}{15}{\milli\meter}$ stratum while heat trickles downward only via slow molecular conduction ($\alpha_w = \SI{1.45e-7}{\square\meter\per\second}$, diffusion depth $\delta_{\mathrm{th}} = \sqrt{4\alpha_w t} \approx \SI{18.7}{\milli\meter}$ after $\SI{10}{\minute}$) \src{S12}.
  \item \textbf{Uncovered Surface Evaporative Skin Cooling ($z = \SIrange{0}{1.5}{\milli\meter}$):}
  The instant the cup is uncovered at the end of Heating Time (\texttt{Rule 5.II.g}), evaporative latent heat loss ($\dot{Q}_{\mathrm{evap}} \approx \SIrange{2.2}{4.4}{\watt}$, \sref{sec:15.4}) cools the uppermost $\SI{1}{\milli\meter}$ free-surface skin by $\SIrange{0.4}{0.9}{\celsius}$ within $\SIrange{15}{30}{\second}$. Consequently, the \textbf{peak water temperature} sits just below the evaporative surface boundary layer at depth $\mathbf{z_{\mathrm{opt}} \approx \SIrange{3}{6}{\milli\meter}}$ (\fref{fig:15:stratification-profile}), inside the $\SIrange{4}{12}{\milli\meter}$ zone an electronic probe tip reaches when the ES follows the FAQ (10/06/2026).
  \item \textbf{Radial Sidewall Cooling Boundary Layer ($r$-Axis Bias):}
  Along the vertical plastic cup wall ($r = R_{\mathrm{water}} \approx \SI{33.8}{\milli\meter}$), radial heat loss to the ambient room cools a $\delta_T \approx \SIrange{3}{5}{\milli\meter}$ natural-convection boundary layer, creating a downward-sinking wall plume that feeds cool water to the cup bottom while leaving the central core ($r < \SI{12}{\milli\meter}$) $\SIrange{0.6}{1.8}{\celsius}$ warmer than the perimeter wall (\fref{fig:17:probe-depth-bias}).
  \item \textbf{Probe Stem Conduction Error at Ultra-Shallow Immersion ($z < \SI{4}{\milli\meter}$):} \idxtopic{stem conduction error in probe}
  A stainless-steel probe sheath ($k_{\mathrm{ss}} \approx \SI{15}{\watt\per\meter\per\kelvin}$, diameter $d_p = \SI{3.2}{\milli\meter}$) acts as a conductive fin protruding into $\SI{22}{\celsius}$ room air. If immersed only $\SI{2}{\milli\meter}$ deep so part of the internal thermistor/RTD element sits above the waterline, axial stem conduction pulls the sensor reading down by $\SIrange{0.8}{2.5}{\celsius}$. Immersing the tip to $z \approx \SIrange{6}{8}{\milli\meter}$ (about 2--2.5 sheath diameters) keeps the bead submerged and still in the hot top layer; at $\SI{10}{\milli\meter}$ the modeled profile already reads $\SI{1.3}{\celsius}$ lower (\fref{fig:17:probe-depth-bias}).
\end{enumerate}

\begin{center}
\resizebox{0.99\linewidth}{!}{%
\begin{tikzpicture}[font=\scriptsize, >=Latex]
  % Left panel: Cross section of 9 oz cup with stratified zones & probe placements
  \begin{scope}[shift={(0,0)}]
    % Stratified water layers (y from 0 to 3.2 cm representing 0 to 38 mm depth)
    \fill[navy!12] (-2.60,0) -- (2.60,0) -- (2.95,1.4) -- (-2.95,1.4) -- cycle;
    \fill[teal!18] (-2.95,1.4) -- (2.95,1.4) -- (3.22,2.45) -- (-3.22,2.45) -- cycle;
    \fill[amber!25] (-3.22,2.45) -- (3.22,2.45) -- (3.35,3.02) -- (-3.35,3.02) -- cycle;
    \fill[rust!22] (-3.35,3.02) -- (3.35,3.02) -- (3.38,3.20) -- (-3.38,3.20) -- cycle;

    % Optimal ES target box (z = 5-12 mm, r < 10 mm)
    \draw[rust, very thick, fill=amber!40] (-1.0,2.20) rectangle (1.0,2.85);
    \node[font=\scriptsize\bfseries, text=rust, align=right, anchor=east] at (-1.2,2.52) {ES target zone\\[-1pt]$z=\SIrange{4}{12}{\milli\meter}$};

    % Cup walls
    \draw[navy, ultra thick] (-3.75,4.8) -- (-2.60,0) -- (2.60,0) -- (3.75,4.8);
    \draw[navy, thick] (-3.85,4.8) -- (-3.65,4.8) (3.65,4.8) -- (3.85,4.8);
    \draw[teal, very thick] (-3.38,3.20) -- (3.38,3.20);

    % Layer labels
    \node[font=\scriptsize, text=navy, align=center] at (-0.35,0.70) {Lower water, $z=\SIrange{24}{38}{\milli\meter}$\\[-1pt]$T = \SIrange{30.8}{36.5}{\celsius}$};
    \node[font=\scriptsize, text=navy, align=center] at (-0.35,1.80) {Thermocline, $z=\SIrange{10}{24}{\milli\meter}$\\[-1pt]$T = \SIrange{36.5}{48.4}{\celsius}$};

    % Radial wall cooling arrows
    \draw[->, navy, thick, dashed] (-2.95,2.05) -- (-2.70,0.95);
    \draw[->, navy, thick, dashed] (2.95,2.05) -- (2.70,0.95);


    % Probe A: Correct ES Center Top-Layer Probe
    \draw[fill=slate!70, draw=black, thick] (-0.12,2.69) rectangle (0.12,5.35);
    \fill[rust] (0,2.69) circle (2.2pt);
    \node[anchor=south east, font=\scriptsize\bfseries, fill=boxbg, draw=navy, inner sep=1.5pt] at (1.85,5.38) {Probe A (ES): center, $z\approx\SI{6}{\milli\meter}$, $\SI{49.7}{\celsius}$};

    % Probe B: Deep/Wall Misplaced Probe
    \draw[fill=midgray!60, draw=black, densely dashed] (2.05,0.25) rectangle (2.25,5.95);
    \fill[navy] (2.15,0.25) circle (2.0pt);
    \node[anchor=south, font=\scriptsize\bfseries, text=rust, fill=white, draw=rust, inner sep=1.5pt] at (2.15,5.98) {Probe B (at the base): $\SI{30.8}{\celsius}$, $-\SI{18.9}{\celsius}$};

    % Depth scale on left
    \draw[<->, black, semithick] (-4.1,3.20) -- (-4.1,0);
    \draw[black] (-4.25,3.20) -- (-3.95,3.20) node[left=2pt] {$z=\SI{0}{\milli\meter}$};
    \draw[black] (-4.25,2.69) -- (-3.95,2.69) node[left=2pt] {$z=\SI{6}{\milli\meter}$};
    \draw[black] (-4.25,0) -- (-3.95,0) node[left=2pt] {$z=\SI{38}{\milli\meter}$};
  \end{scope}

  % Right panel: Vertical & Radial Temperature Profile Table / Mini-Plot
  \begin{scope}[shift={(5.5,0.2)}]
    \draw[navy, thick, fill=boxbglight] (0,-0.35) rectangle (6.5,5.0);
    \node[anchor=north, font=\scriptsize\bfseries, fill=navy, text=white, minimum width=6.25cm, inner sep=2.5pt] at (3.25,5.0) {Modeled Profile $T(z)$ at Cup Center ($t=\SI{10}{\minute}$, $h=\SI{45}{\centi\meter}$)};
    \node[anchor=north west, font=\scriptsize] at (0.35,4.45) {%
      \setlength{\tabcolsep}{4pt}
      \renewcommand{\arraystretch}{1.12}
      \begin{tabular}{@{}lcc@{}}
      \toprule
      \textbf{Probe tip depth $z$} & \textbf{$T$ at center} & \textbf{Bias vs.\ ES} \\
      \midrule
      $\SI{0.5}{\milli\meter}$ (skin; IR gun) & $\SI{49.4}{\celsius}$ & $-\SI{0.3}{\celsius}$ \\
      $\SI{3}{\milli\meter}$ (peak) & $\SI{49.9}{\celsius}$ & $+\SI{0.2}{\celsius}$ \\
      $\mathbf{\SI{6}{\milli\meter}}$ \textbf{(ES reading)} & $\mathbf{\SI{49.7}{\celsius}}$ & $\mathbf{\SI{0.0}{\celsius}}$ \textbf{(ref)} \\
      $\SI{10}{\milli\meter}$ & $\SI{48.4}{\celsius}$ & $-\SI{1.3}{\celsius}$ \\
      $\SI{16}{\milli\meter}$ & $\SI{43.2}{\celsius}$ & $-\SI{6.5}{\celsius}$ \\
      $\SI{24}{\milli\meter}$ & $\SI{36.5}{\celsius}$ & $-\SI{13.2}{\celsius}$ \\
      $\SI{32}{\milli\meter}$ & $\SI{32.0}{\celsius}$ & $-\SI{17.7}{\celsius}$ \\
      $\SI{38}{\milli\meter}$ (base) & $\SI{30.8}{\celsius}$ & $\mathbf{-\SI{18.9}{\celsius}}$ \\
      \midrule
      Swirled (mixed; not allowed) & $\SI{39.8}{\celsius}$ & $\mathbf{-\SI{9.9}{\celsius}}$ \\
      \bottomrule
      \end{tabular}%
    };
    \node[anchor=south, font=\scriptsize\bfseries, text=rust, align=center] at (3.25,-0.30) {Modeled for the reference collector.\\[-1pt]Replace with our measured runs.};
  \end{scope}
\end{tikzpicture}%
}
\captionof{figure}{Modeled probe-depth bias in the reference collector's unstirred, top-heated $\SI{9}{oz}$ cup ($\SI{100}{\milli\liter}$ water, $t=\SI{10}{\minute}$, $h=\SI{45}{\centi\meter}$, $\SI{125}{\watt}$; same model as \fref{fig:15:stratification-profile} and Ex~\ref{ex:15.4}). The ES reads the top layer near the center (FAQ 10/06/2026): $\SI{49.7}{\celsius}$. A probe pushed to the base reads $\SI{18.9}{\celsius}$ low; swirled water would read the mixed $\SI{39.8}{\celsius}$, but swirling is not allowed (FAQs 09/30, 10/01 and 10/06/2026). Dashed arrows: the cooler wall plume (about $\SI{1.4}{\celsius}$ below the center at the same depth) \src{S02,S05,S19}. \textbf{How to read this:} find your probe depth in the table; the last column is how far that reading falls below the ES's.}\label{fig:17:probe-depth-bias}
\idxfig{ES Thermometer Immersion Depth and Radial Placement Bias in Stratified Cup}
\end{center}

% ============================================================================
\section{Fixed-Point Thermometer Calibration: Ice Bath ($\SI{0.00}{\celsius}$) \& Barometric Steam Point}\label{sec:17.5}
\idxtopic{thermometer calibration!ice point and boiling point}\idxformula{barometric boiling point T_boil = 100.00 + 0.03686*(P_mmHg - 760) - 0.000020*(P_mmHg - 760)^2}

To verify or calibrate team thermometers before building the 36-run DOE calibration table (\sref{sec:16.2}), perform a two-point fixed-temperature calibration against the \textbf{ice point ($\SI{0.00}{\celsius} = \SI{273.15}{\kelvin}$)} and the \textbf{barometric-corrected steam boiling point ($T_{\mathrm{boil}}(P_{\mathrm{atm}})$)} \src{S09,S21,S22}:
\begin{enumerate}
  \item \textbf{Ice-Point Reference Bath ($T_{\mathrm{ice}} = \SI{0.00}{\celsius}$ at $\SI{1}{\text{atm}}$):} \idxerr{ice floating on 4 C warm bottom water in calibration bath}
  Pack a clean Dewar or insulated foam cup \textbf{all the way to the bottom} with finely shaved or crushed ice made from distilled water, then add just enough chilled distilled water to fill the pore spaces between ice crystals ($\SIrange{75}{80}{\percent}$ solid ice, $\SIrange{20}{25}{\percent}$ liquid water; slush looks wet and translucent, with \textbf{zero free-floating ice gap at the bottom}).
  \emph{Why bottom-packed ice is critical:} Because liquid water reaches maximum density $\rho_{\max} = \SI{999.97}{\kilogram\per\cubic\meter}$ at $+\SI{3.984}{\celsius}$ (\sref{sec:03}) \src{S09}, if ice cubes merely float at the top of a cup of water, dense $+\SI{4}{\celsius}$ water sinks to the bottom, causing a deep probe to read $+\SIrange{1.5}{4.0}{\celsius}$ instead of $\SI{0.00}{\celsius}$! Wait $\SI{3}{\minute}$, insert the probe into the central slush zone without touching the container walls, wait $\ge 5\tau_p$, and record $T_{\mathrm{read,ice}}$. The zero-offset error is $\delta_0 = T_{\mathrm{read,ice}} - \SI{0.00}{\celsius}$. \emph{(Note: Dissolved air depresses the freezing point of air-saturated water by only $\SI{-0.0024}{\celsius}$ relative to the $\SI{0.0100}{\celsius}$ triple point, and barometric pressure shifts the ice point by $\frac{dT_m}{dP} = \SI{-0.00743}{\celsius\per\text{atm}}$, which is negligible on a $\SI{0.1}{\celsius}$ probe \src{S11}).}
  \item \textbf{Altitude / Barometric-Pressure Corrected Steam Boiling Point ($T_{\mathrm{boil}}(P_{\mathrm{atm}})$):}
  Water boils at exactly $\SI{100.00}{\celsius}$ ($T_{90} = \SI{99.974}{\celsius}$ on ITS-90 \src{S09}) \textbf{only} at standard sea-level pressure $P_0 = \SI{101.325}{\kilo\pascal} = \SI{760.00}{\milli\meter Hg}$. Because $P_{\mathrm{atm}}$ drops with elevation $z_{\mathrm{alt}}$ $[\si{\meter}]$ as $P_{\mathrm{atm}}(z_{\mathrm{alt}}) \approx \SI{101.325}{\kilo\pascal}\left(1 - \num{2.2558e-5}\,z_{\mathrm{alt}}\right)^{5.2559}$ ($\approx -\SI{1.2}{\kilo\pascal}$ or $-\SI{9.0}{\milli\meter Hg}$ per $+\SI{100}{\meter}$ elevation) and fluctuates by $\pm\SI{2.5}{\kilo\pascal}$ with weather fronts, calibrating a thermometer in boiling water without correcting for local barometric pressure introduces a $\SIrange{0.5}{3.0}{\celsius}$ systematic error! From the NIST IAPWS-95 saturation curve around $\SI{100}{\celsius}$ ($\frac{dT_{\mathrm{sat}}}{dP} = \frac{T v_{fg}}{h_{fg}} = \SI{0.2765}{\kelvin\per\kilo\pascal} = \SI{0.03686}{\kelvin\per\milli\meter Hg}$) \src{S09}:
  \begin{align}\label{eq:17:tboil}
    T_{\mathrm{boil}}(P_{\mathrm{mmHg}}) &= \SI{100.00}{\celsius} + \num{0.03686}\left(P_{\mathrm{mmHg}} - 760.0\right) - \num{0.000020}\left(P_{\mathrm{mmHg}} - 760.0\right)^2 \\[2pt]
    T_{\mathrm{boil}}(P_{\mathrm{kPa}})  &= \SI{100.00}{\celsius} + \num{0.2765}\left(P_{\mathrm{kPa}} - 101.325\right) - \num{0.00112}\left(P_{\mathrm{kPa}} - 101.325\right)^2
  \end{align}
  \emph{(Suspend the probe in the condensing saturated steam just above the boiling distilled water surface or in a vigorously rolling bath with boiling chips to avoid $+\SIrange{0.2}{0.6}{\celsius}$ liquid superheating at the heated vessel bottom).}
  \item \textbf{Two-Point Linear Calibration Correction Equation:} \idxformula{two-point thermometer calibration T_true = T_boil*(T_raw - T_read,ice)/(T_read,boil - T_read,ice)}
  Given raw thermometer reading $T_{\mathrm{raw}}$, ice-bath reading $T_{\mathrm{read,ice}}$ (at true $\SI{0.00}{\celsius}$), and steam-point reading $T_{\mathrm{read,boil}}$ (at true $T_{\mathrm{boil}}(P_{\mathrm{atm}})$):
  \begin{equation}\label{eq:17:2ptcal}
    \boxed{T_{\mathrm{true}} = T_{\mathrm{boil}}(P_{\mathrm{atm}}) \times \frac{T_{\mathrm{raw}} - T_{\mathrm{read,ice}}}{T_{\mathrm{read,boil}} - T_{\mathrm{read,ice}}}}
  \end{equation}
\end{enumerate}

\begin{workedexample}[5$\tau_p$ Premature-Reading Audit, Probe Heat Drain \& Two-Point Barometric Calibration]{Rule 4.f \& 5.II.g}{ex:17.4}
(a)~A stainless-sheathed digital probe ($\tau_p = \SI{3.20}{\second}$, effective tip heat capacity $C_p = \SI{0.450}{\joule\per\kelvin}$, initial room temperature $T_{p,0} = \SI{21.50}{\celsius}$) is inserted into the top stratified water layer ($C_{\mathrm{top}} = \SI{215.0}{\joule\per\kelvin}$, initial temperature $T_{w,0} = \SI{49.70}{\celsius}$). Compute the true equilibrium temperature $T_{\mathrm{eq}}$ after probe thermal drain, the premature reading if recorded at $t = \SI{6.40}{\second}$ ($2\tau_p$), and the settled reading at $t = 5\tau_p = \SI{16.0}{\second}$.
(b)~During pre-tournament calibration on a day when local barometric pressure is $P_{\mathrm{atm}} = \SI{742.0}{\milli\meter Hg}$, the probe reads $T_{\mathrm{read,ice}} = -\SI{0.30}{\celsius}$ in an ice bath and $T_{\mathrm{read,boil}} = \SI{98.90}{\celsius}$ in saturated steam. If the probe reads $T_{\mathrm{raw}} = \SI{44.20}{\celsius}$ in a trial cup, what is the true calibrated water temperature $T_{\mathrm{true}}$?
\par\smallskip\noindent
\textbf{Solution.}
(a)~By \eref{eq:17:probedrain}, calorimetric equilibration between the $\SI{21.50}{\celsius}$ probe tip and the $\SI{49.70}{\celsius}$ top layer yields:
$$T_{\mathrm{eq}} = \frac{215.0(49.70) + 0.450(21.50)}{215.0 + 0.450} = \frac{10695.175}{215.450} = \mathbf{\SI{49.641}{\celsius}} \quad (\Delta T_{\mathrm{drain}} = -\SI{0.059}{\celsius})$$
Using \eref{eq:17:probestep} with $\Delta T_{\mathrm{step}} = 49.6411 - 21.50 = \SI{28.1411}{\celsius}$:
$$T_p(2\tau_p = \SI{6.40}{\second}) = 49.6411 - 28.1411\,e^{-2} = \mathbf{\SI{45.83}{\celsius}} \quad (\text{premature reading loses }\mathbf{\SI{3.81}{\celsius}}\text{!})$$
$$T_p(5\tau_p = \SI{16.0}{\second}) = 49.6411 - 28.1411\,e^{-5} = \mathbf{\SI{49.45}{\celsius}} \quad (\text{within }\SI{0.19}{\celsius}\text{ of }T_{\mathrm{eq}})$$
(b)~At $P_{\mathrm{atm}} = \SI{742.0}{\milli\meter Hg}$ ($\Delta P = -\SI{18.0}{\milli\meter Hg}$), \eref{eq:17:tboil} gives true steam temperature:
$$T_{\mathrm{boil}}(\SI{742.0}{\milli\meter Hg}) = 100.00 + \num{0.03686}(-18.0) - \num{0.000020}(-18.0)^2 = 100.00 - 0.66348 - 0.00648 = \mathbf{\SI{99.330}{\celsius}}$$
Applying the two-point linear calibration formula \eref{eq:17:2ptcal}:
$$T_{\mathrm{true}} = \SI{99.330}{\celsius} \times \frac{44.20 - (-0.30)}{98.90 - (-0.30)} = 99.330 \times \frac{44.50}{99.20} = \mathbf{\SI{44.56}{\celsius}}$$
\verifiedby{Direct linear slope/offset verification $T_{\mathrm{true}} = 1.00131(44.20 + 0.30) = \SI{44.5583}{\celsius}$ and Python check \texttt{sec17:ex17.4}.}
\end{workedexample}

% ============================================================================
\section{JCGM GUM / NIST TN~1297 Uncertainty Propagation Framework \src{S21}}\label{sec:17.6}
\idxtopic{uncertainty propagation (JCGM GUM / NIST TN 1297)}\idxtopic{Type A vs Type B standard uncertainty}\idxformula{GUM propagation law u_c^2(y) = sum (df/dx_i)^2 * u^2(x_i)}

Under the international BIPM/JCGM \emph{Guide to the Expression of Uncertainty in Measurement} (GUM, JCGM~100:2008) and NIST Technical Note~1297 \src{S21}, every measurement uncertainty is expressed as a \textbf{standard uncertainty} $u(x_i)$ (equivalent to one standard deviation, $1\sigma$) classified by evaluation method:
\begin{enumerate}
  \item \textbf{Type~A Standard Uncertainty ($u_A$ --- Statistical Evaluation of Repeated Trials):}
  From $n$ independent repeated measurements $x_1, x_2, \dots, x_n$ with sample mean $\bar{x} = \frac{1}{n}\sum_{k=1}^n x_k$ and sample standard deviation $s = \sqrt{\frac{1}{n-1}\sum_{k=1}^n (x_k - \bar{x})^2}$ (with $\nu = n - 1$ degrees of freedom):
  \begin{equation}\label{eq:17:typeA}
    u_A(\bar{x}) = \frac{s}{\sqrt{n}} = \sqrt{\frac{\sum_{k=1}^n (x_k - \bar{x})^2}{n(n - 1)}}
  \end{equation}
  \item \textbf{Type~B Standard Uncertainty ($u_B$ --- Non-Statistical Prior / Specification Bounds):}
  \begin{itemize}
    \item \emph{Manufacturer Tolerance $\pm a$ (Rectangular / Uniform Distribution on $[\bar{x}-a, \bar{x}+a]$):}
    \begin{equation}\label{eq:17:typeBtol}
      u_{B,\mathrm{tol}} = \frac{a}{\sqrt{3}} \approx 0.57735\,a
    \end{equation}
    \item \emph{Digital Display Resolution $\Delta_{\mathrm{res}}$ (Uniform Rounding Within $\pm a = \pm\Delta_{\mathrm{res}}/2$):}
    \begin{equation}\label{eq:17:typeBres}
      u_{B,\mathrm{res}} = \frac{\Delta_{\mathrm{res}} / 2}{\sqrt{3}} = \frac{\Delta_{\mathrm{res}}}{2\sqrt{3}} \approx 0.28868\,\Delta_{\mathrm{res}} \quad (\text{for }\Delta_{\mathrm{res}} = \SI{0.1}{\celsius},\ u_{B,\mathrm{res}} = \SI{0.0289}{\celsius})
    \end{equation}
    \item \emph{Analog Meniscus / Scale Reading Bound $\pm a$ (Triangular Distribution Peaked at Center):} $u_{B,\mathrm{tri}} = \frac{a}{\sqrt{6}} \approx 0.40825\,a$.
  \end{itemize}
  \item \textbf{Combined Standard Uncertainty $u_c(y)$ \& Expanded Uncertainty $U = k\,u_c(y)$:}
  For a single measurand $x$ with independent Type~A and Type~B components, $u(x) = \sqrt{u_A^2 + u_{B,\mathrm{tol}}^2 + u_{B,\mathrm{res}}^2}$. For a derived quantity $y = f(x_1, x_2, \dots, x_N)$, the first-order Taylor expansion (Law of Propagation of Uncertainty) with sensitivity coefficients $c_i = \frac{\partial f}{\partial x_i}$ and covariance $u(x_i, x_j) = r_{ij}u(x_i)u(x_j)$ gives \src{S21}:
  \begin{equation}\label{eq:17:gum}
    u_c^2(y) = \sum_{i=1}^N \left(\frac{\partial f}{\partial x_i}\right)^2 u^2(x_i) + 2\sum_{i=1}^{N-1}\sum_{j=i+1}^N \frac{\partial f}{\partial x_i}\frac{\partial f}{\partial x_j}\,r_{ij}u(x_i)u(x_j), \qquad U = k\,u_c(y)
  \end{equation}
  where coverage factor $k = 2$ defines an approximately $\SI{95}{\percent}$ confidence interval $y \pm U$ (or $k = t_{0.95,\nu_{\mathrm{eff}}}$ from the Student-$t$ distribution when effective degrees of freedom $\nu_{\mathrm{eff}}$ from the Welch--Satterthwaite formula $\nu_{\mathrm{eff}} = u_c^4(y)/\sum \frac{[c_i u(x_i)]^4}{\nu_i}$ are small, e.g., $\nu < 10$).
\end{enumerate}

\begin{table}[h!]
\centering
\captionof{table}{Algebra-Only JCGM GUM Uncertainty Propagation Rules for Science Olympiad Thermodynamics (`Rule 5.d.iv` Compliant) \src{S06,S21}.}\label{tab:17:gumrules}
\footnotesize
\setlength{\tabcolsep}{3.5pt}
\renewcommand{\arraystretch}{1.22}
\begin{tabularx}{\linewidth}{@{}>{\raggedright\arraybackslash}p{0.21\linewidth} >{\raggedright\arraybackslash}p{0.24\linewidth} >{\raggedright\arraybackslash}p{0.31\linewidth} >{\raggedright\arraybackslash}X@{}}
\toprule
\hdr{Functional Relationship} & \hdr{Equation $y = f(x_1, x_2, \dots)$} & \hdr{Ready-to-Use Algebraic Uncertainty Rule} & \hdr{Thermodynamics Competition Application} \\
\midrule
\textbf{1.~Linear Sum / Difference} (Uncorrelated, $r_{ij}=0$) &
$y = a x_1 + b x_2 - c x_3$ &
$u_c(y) = \sqrt{a^2 u^2(x_1) + b^2 u^2(x_2) + c^2 u^2(x_3)}$ \newline \emph{(Absolute uncertainties add in quadrature)} &
$\Delta T = T_2 - T_1$ with two \emph{different} sensors: $u(\Delta T) = \sqrt{u^2(T_2) + u^2(T_1)}$. \\
\rowcolor{rowalt}
\textbf{2.~Correlated Difference} (\textbf{Rule \texttt{4.f} Same Thermometer!}) \idxtopic{Rule 4.f!GUM correlated uncertainty} &
$\mathrm{TG} = T_f - T_s$ \newline ($r_{\mathrm{sys}} = +1$ for shared bias $u_{B,\mathrm{sys}}$) &
$u_c(\mathrm{TG}) = \sqrt{u_{\mathrm{rand}}^2(T_f) + u_{\mathrm{rand}}^2(T_s)}$ \newline where $u_{\mathrm{rand}}^2 = u_A^2 + u_{B,\mathrm{res}}^2$ \newline \emph{(Shared systematic tolerance $u_{B,\mathrm{sys}}$ cancels!)} &
Under Rule~\texttt{4.f}, the same thermometer measures $T_s$ and $T_f$: $u_{B,\mathrm{sys}}^2 + u_{B,\mathrm{sys}}^2 - 2(+1)u_{B,\mathrm{sys}}^2 = 0$! \\
\textbf{3.~Product \& Quotient} &
$y = \dfrac{k\,x_1\,x_2}{x_3}$ &
$\dfrac{u_c(y)}{|y|} = \sqrt{\left(\dfrac{u(x_1)}{x_1}\right)^2 + \left(\dfrac{u(x_2)}{x_2}\right)^2 + \left(\dfrac{u(x_3)}{x_3}\right)^2}$ \newline \emph{(Fractional uncertainties add in quadrature)} &
Calorimetry $Q = m c_p \Delta T$, ideal gas $P = nRT/V$, or thermal conductivity $k = \dfrac{\dot{Q} L}{A \Delta T}$. \\
\rowcolor{rowalt}
\textbf{4.~Power / Monomial Law} &
$y = k\,x_1^a\,x_2^b\,x_3^{-c}$ &
$\dfrac{u_c(y)}{|y|} = \sqrt{\left(a\dfrac{u(x_1)}{x_1}\right)^2 + \left(b\dfrac{u(x_2)}{x_2}\right)^2 + \left(c\dfrac{u(x_3)}{x_3}\right)^2}$ &
Stefan--Boltzmann $\dot{Q} = \varepsilon\sigma A T^4 \implies \dfrac{u(\dot{Q})}{\dot{Q}} = \sqrt{\dots + \left(\mathbf{4}\dfrac{u(T)}{T}\right)^2}$; pipe area $A = \frac{\pi}{4}D^2 \implies 2\frac{u(D)}{D}$. \\
\textbf{5.~Natural Logarithm \& Exponential} &
$y = a\ln(x)$ \newline or $z = a e^{b x}$ &
$u_c(y) = |a|\dfrac{u(x)}{|x|}$ \newline $\dfrac{u_c(z)}{|z|} = |b|\,u(x)$ &
Semi-log cooling plot $y = \ln(T - T_{\mathrm{ss}})$, entropy $\Delta S = mc\ln(T_2/T_1)$, or Boltzmann factor $e^{-E/(k_B T)}$. \\
\rowcolor{rowalt}
\textbf{6.~OLS Linear Regression Slope $m$ \& Intercept $b$} &
$y = m x + b$ across $N$ points $(x_i, y_i)$ &
$s_{y|x} = \sqrt{\dfrac{S_{yy} - m S_{xy}}{N - 2}}$,\quad $u(m) = \dfrac{s_{y|x}}{\sqrt{S_{xx}}}$, \newline $u(b) = s_{y|x}\sqrt{\dfrac{1}{N} + \dfrac{\bar{x}^2}{S_{xx}}}$ &
Time constant $\tau = -1/m \implies u(\tau) = \dfrac{u(m)}{m^2} = \tau^2 u(m)$; Clausius--Clapeyron $\Delta H_{\mathrm{vap}} = -R m \implies u(\Delta H) = R\,u(m)$. \\
\bottomrule
\end{tabularx}
\end{table}

\begin{workedexample}[Complete JCGM GUM Uncertainty Budget for \texttt{Rule 4.f} $\mathrm{TG}$ \& Specific-Heat Calorimetry]{Rules 4.f \& 5.I.b}{ex:17.5}
(a)~An ES digital probe has manufacturer tolerance $\pm a_{\mathrm{sys}} = \pm\SI{0.40}{\celsius}$, display resolution $\Delta_{\mathrm{res}} = \SI{0.10}{\celsius}$, and Type~A random repeatability $u_A = \SI{0.050}{\celsius}$. Calculate the combined standard uncertainty $u(T)$ of a single temperature reading, and compare $u_c(\mathrm{TG})$ for $\mathrm{TG} = T_f - T_{\mathrm{start}}$ when measured with two \emph{different} probes ($r=0$) versus the \emph{same} probe under Rule~\texttt{4.f} ($r_{\mathrm{sys}} = +1$).
(b)~In a calorimetry lab using that same thermometer, a metal sample of mass $m_x = \SI{85.00 \pm 0.05}{\gram}$ cooled by $\Delta T_x = \SI{52.60}{\celsius}$ warms $m_w = \SI{100.00 \pm 0.15}{\gram}$ of water ($c_w = \SI{4.184}{\joule\per\gram\per\kelvin}$, exact) by $\Delta T_w = \SI{5.40}{\celsius}$ (where $u(\Delta T_w) = u(\Delta T_x) = u_{\mathrm{corr}}(\mathrm{TG}) = \SI{0.08165}{\celsius}$). Compute $c_x$, its combined standard uncertainty $u_c(c_x)$, and the expanded uncertainty $U(c_x)$ ($k=2$).
\par\smallskip\noindent
\textbf{Solution.}
(a)~From \eref{eq:17:typeBtol}--\eref{eq:17:typeBres}: $u_{B,\mathrm{sys}} = \frac{0.40}{\sqrt{3}} = \SI{0.23094}{\celsius}$ and $u_{B,\mathrm{res}} = \frac{0.10}{2\sqrt{3}} = \SI{0.02887}{\celsius}$. For a single reading:
$$u(T) = \sqrt{u_{B,\mathrm{sys}}^2 + u_{B,\mathrm{res}}^2 + u_A^2} = \sqrt{0.23094^2 + 0.02887^2 + 0.050^2} = \mathbf{\SI{0.2380}{\celsius}}$$
If two independent thermometers were used ($r = 0$), $u_{\mathrm{uncorr}}(\mathrm{TG}) = \sqrt{2}\,u(T) = \mathbf{\SI{0.3367}{\celsius}}$. Under Rule~\texttt{4.f} (same thermometer, $r_{\mathrm{sys}} = +1$), the correlated systematic variance cancels ($u_{B,\mathrm{sys}}^2 + u_{B,\mathrm{sys}}^2 - 2u_{B,\mathrm{sys}}^2 = 0$), leaving only random and resolution terms:
$$u_{\mathrm{corr}}(\mathrm{TG}) = \sqrt{2\left(u_{B,\mathrm{res}}^2 + u_A^2\right)} = \sqrt{2\left(0.02887^2 + 0.050^2\right)} = \mathbf{\SI{0.08165}{\celsius}} \quad (\mathbf{4.12\times\text{\textbf{ lower uncertainty!}}})$$
(b)~Specific heat of the metal: $c_x = \frac{m_w c_w \Delta T_w}{m_x \Delta T_x} = \frac{(100.00)(4.184)(5.40)}{(85.00)(52.60)} = \mathbf{\SI{0.50534}{\joule\per\gram\per\kelvin}}$. By Rule~3 in Table~\ref{tab:17:gumrules}:
$$\frac{u_c(c_x)}{c_x} = \sqrt{\left(\frac{0.15}{100.00}\right)^2 + \left(\frac{0.05}{85.00}\right)^2 + \left(\frac{0.08165}{5.40}\right)^2 + \left(\frac{0.08165}{52.60}\right)^2} = \mathbf{0.01528}\ (\SI{1.53}{\percent})$$
\emph{(Notice that $\frac{u(\Delta T_w)}{\Delta T_w} = \SI{1.51}{\percent}$ contributes $>97\%$ of the total variance!)} Thus $u_c(c_x) = (0.50534)(0.01528) = \mathbf{\SI{0.00772}{\joule\per\gram\per\kelvin}}$ and expanded uncertainty ($k=2$) is $U = 2 u_c(c_x) = \mathbf{\SI{0.0154}{\joule\per\gram\per\kelvin}}$, reported as $c_x = \mathbf{\SI{0.505 \pm 0.015}{\joule\per\gram\per\kelvin}}$.
\verifiedby{Monte Carlo / finite-difference perturbation of $c_x(m_w, m_x, \Delta T_w, \Delta T_x)$ and Python check \texttt{sec17:ex17.5}.}
\end{workedexample}

\begin{workedexample}[Thermocouple Cold-Junction Compensation (CJC) \& Stefan--Boltzmann $T^4$ Power-Law Uncertainty]{Rules 4.f, 5.c.vi}{ex:17.6}
(a)~A \texttt{Type T} thermocouple has its sensing junction in warm water at $T_{\mathrm{sense}}$ and its voltmeter terminal block (reference junction) at ambient room temperature $T_{\mathrm{ref}} = \SI{20.00}{\celsius}$. The uncompensated millivoltmeter reads $\Delta V(T_{\mathrm{ref}}, T_{\mathrm{sense}}) = \SI{+0.934}{\milli\volt}$. Using Table~\ref{tab:17:sensortable}, find the total ice-point referenced EMF $E(\SI{0}{\celsius}, T_{\mathrm{sense}})$ and determine $T_{\mathrm{sense}}$ by linear interpolation between $\SI{40.0}{\celsius}$ and $\SI{50.0}{\celsius}$.
(b)~A blackbody cavity ($\varepsilon = 1.000$ exact) of aperture area $A = \SI{50.00 \pm 0.25}{\square\centi\meter}$ operates at $T = \SI{500.0 \pm 2.5}{\kelvin}$. Calculate the emitted radiant power $\dot{Q} = \sigma A T^4$ and its combined standard uncertainty $u_c(\dot{Q})$.
\par\smallskip\noindent
\textbf{Solution.}
(a)~From Table~\ref{tab:17:sensortable}, a \texttt{Type T} junction at $T_{\mathrm{ref}} = \SI{20.00}{\celsius}$ has $E(\SI{0}{\celsius}, \SI{20.00}{\celsius}) = \SI{+0.790}{\milli\volt}$. By the Law of Intermediate Temperatures:
$$E(\SI{0}{\celsius}, T_{\mathrm{sense}}) = E(\SI{0}{\celsius}, \SI{20.00}{\celsius}) + \Delta V(\SI{20.00}{\celsius}, T_{\mathrm{sense}}) = 0.790 + 0.934 = \mathbf{\SI{+1.724}{\milli\volt}}$$
Interpolating in Table~\ref{tab:17:sensortable} between $\SI{40.0}{\celsius}$ ($\SI{1.612}{\milli\volt}$) and $\SI{50.0}{\celsius}$ ($\SI{2.036}{\milli\volt}$):
$$T_{\mathrm{sense}} = 40.0 + (50.0 - 40.0)\frac{1.724 - 1.612}{2.036 - 1.612} = 40.0 + 10.0\left(\frac{0.112}{0.424}\right) = \mathbf{\SI{42.64}{\celsius}}$$
(b)~Emitted radiant power ($\sigma = \SI{5.670374e-8}{\watt\per\square\meter\per\kelvin\tothe{4}}$, $A = \SI{5.000e-3}{\square\meter}$):
$$\dot{Q} = \sigma A T^4 = (\num{5.670374e-8})(\num{5.000e-3})(500.0)^4 = \mathbf{\SI{17.720}{\watt}}$$
By the power-law GUM formula (Rule~4 in Table~\ref{tab:17:gumrules}), the exponent $4$ on $T$ multiplies its relative uncertainty by $4$:
$$\frac{u_c(\dot{Q})}{\dot{Q}} = \sqrt{\left(\frac{u(A)}{A}\right)^2 + \left(4\,\frac{u(T)}{T}\right)^2} = \sqrt{\left(\frac{0.25}{50.00}\right)^2 + \left(4\times\frac{2.5}{500.0}\right)^2} = \sqrt{0.0050^2 + 0.0200^2} = \mathbf{0.02062}\ (\SI{2.06}{\percent})$$
Thus $u_c(\dot{Q}) = (\SI{17.720}{\watt})(0.020616) = \mathbf{\SI{0.365}{\watt}}$, so $\dot{Q} = \mathbf{\SI{17.72 \pm 0.37}{\watt}}$ ($1\sigma$).
\verifiedby{Exact differential $\delta\dot{Q} = \sqrt{(\sigma T^4 u_A)^2 + (4\sigma A T^3 u_T)^2} = \sqrt{0.0886^2 + 0.3544^2} = \SI{0.3653}{\watt}$ and Python check \texttt{sec17:ex17.6}.}
\end{workedexample}

% ============================================================================
\section{Step-by-Step Classroom-Safe Calorimetry \& 2027 Collector Validation Protocols}\label{sec:17.7}
\idxtopic{laboratory protocols!calorimetry and collector validation}\idxtopic{safety checklist!burn, electrical, spill, glass}

\begin{devicebox}[Protocol A (Calorimeter Heat Capacity $C_{\mathrm{cal}}$) \& Protocol B (2027 Heat-Lamp Collector DOE Validation)]
\begin{twocol}[0.485]
\textbf{Protocol A: Classroom Calorimeter Calibration \& Specific Heat Determination (\texttt{Rule 5.c.iii})}
\begin{enumerate}
  \item \textbf{Dry Tare \& Cold Baseline:} Weigh dry inner calorimeter cup ($m_{\mathrm{cup}}$) on $\SI{0.01}{\gram}$ balance. Add $\SI{100.0}{\milli\liter}$ room-temperature water, re-weigh ($m_c = m_{\mathrm{total}} - m_{\mathrm{cup}}$), insert calibrated probe, stir gently, wait $\ge 5\tau_p$, and log $T_c$ every $\SI{15}{\second}$ for $\SI{2}{\minute}$ to confirm zero drift.
  \item \textbf{Determine Calorimeter Constant $C_{\mathrm{cal}}$ $[\si{\joule\per\kelvin}]$:} Heat a separate $\SI{100.0}{\milli\liter}$ water sample to $T_h \approx \SI{55}{\celsius}$ (safe below scalding threshold). Record $T_h$ ($\ge 5\tau_p$) and exact mass $m_h$, pour rapidly into cold calorimeter, cover lid, and log $T(t)$ every $\SI{5}{\second}$ through peak $T_{\mathrm{eq}}$ and for $\SI{3}{\minute}$ of post-peak cooling. Extrapolate the post-peak linear cooling line back to mixing time $t_{\mathrm{mix}}$ (Regnault--Pfaundler cooling correction, \sref{sec:04}) to get true $T_{\mathrm{eq}}$, then compute:
  $$C_{\mathrm{cal}} = \frac{m_h c_w (T_h - T_{\mathrm{eq}}) - m_c c_w (T_{\mathrm{eq}} - T_c)}{T_{\mathrm{eq}} - T_c}$$
  \item \textbf{Unknown Solid Specific Heat $c_x$:} Equilibrate dry metal cylinder ($m_x \approx \SI{80}{\gram}$) in a boiling steam jacket ($T_{x,i} = T_{\mathrm{boil}}(P_{\mathrm{atm}})$ for $\ge\SI{10}{\minute}$), transfer rapidly ($<\SI{2}{\second}$ in air) into fresh calorimeter water ($m_w, T_{w,i}$) without splashing, record $T_{\mathrm{eq}}$, and solve:
  $$c_x = \frac{\left(m_w c_w + C_{\mathrm{cal}}\right)(T_{\mathrm{eq}} - T_{w,i})}{m_x\left(T_{x,i} - T_{\mathrm{eq}}\right)}$$
\end{enumerate}
\nextcol[0.485]
\textbf{Protocol B: 2027 Heat-Lamp Collector Calibration \& Stratification Validation (\texttt{Rule 5.II})}
\begin{enumerate}
  \item \textbf{Station Geometry \& Lamp Pre-Warm:} Mount clamp lamp with $\SI{125}{\watt}$ (or $\SI{100}{\watt}$) clear untinted incandescent heat-lamp bulb facing straight down (\texttt{Rule 4.a--d}). Set vertical height $h \in \SIrange{40.0}{65.0}{\centi\meter}$ from table surface to bottommost tip of bulb (\texttt{Rule 4.c}) using a rigid meterstick and plumb line. Turn lamp ON for $\ge\SI{3}{\minute}$ off to the side so filament and glass envelope reach thermal steady state.
  \item \textbf{3-Min Device Setup \& 1-Min Water Setup Simulation:} Start master timer. Configure collector pedestal slot for height $h$ (\texttt{Rule 5.II.b}). Dispense $\SI{99.67 \pm 0.05}{\gram}$ ($\SI{100.00}{\milli\liter}$ at $\SI{22}{\celsius}$) water from source bath into a fresh $\SI{9}{oz}$ clear cup, measure $T_{\mathrm{start}}$ at $z = \SI{8}{\milli\meter}$ after $5\tau_p$, place non-attached clear BoPET film cap on cup rim (FAQ 10/21/2026 on items on the cup), seat cup into collector cradle, and center under lamp at $t = 0$.
  \item \textbf{Heating Window \& Unswirled Top-Layer Extraction:} Heat for exact target duration $t \in \{10, 12, 15\}\,\si{\minute}$. At lamp turn-off, lift upper reflector hood straight up (removing the clear film cap cleanly) and lift the $\SI{9}{oz}$ cup vertically onto the Testing Surface in $<\SI{8}{\second}$ \textbf{with zero tilting, swirling, or vibration} (FAQs 09/30, 10/01 and 10/06/2026 on stirring and swirling).
  \item \textbf{FAQ 10/06/2026 on probe depth Probe Reading \& Cool-Down Return:} Lower probe clamp into top water layer ($z = \SI{8 \pm 2}{\milli\meter}$, $r = \SI{0}{\milli\meter}$), wait $5\tau_p$ ($\SI{10}{\second}$), and log $T_f$ in Appendix~D test log. Cool collector back to $T_{\mathrm{room}}$ with a fan before the next trial!
\end{enumerate}
\end{twocol}
\end{devicebox}

\begin{rulebox}[Mandatory Laboratory Safety \& Impound Compliance Checklist (Thermal, Electrical, Glass \& Chemical)]
\begin{itemize}
  \item \textbf{Thermal Burn \& Fire Prevention:} A $\SI{125}{\watt}$ heat lamp glass face reaches $\SIrange{180}{230}{\celsius}$; never touch the bulb or metal clamp shroud until cooled for $\ge\SI{15}{\minute}$. Keep all paper/cardboard/balsa collector surfaces $\ge\SI{10}{\centi\meter}$ below the bulb and faced with reflective aluminum foil ($\rho_{\mathrm{NIR}} = 0.96$) so structural temperatures stay $<\SI{45}{\celsius}$ (well below the $\SI{218}{\celsius}$ paper ignition point and $\SI{74}{\celsius}$ \texttt{\#1 PET} cup softening point).
  \item \textbf{Zero-Tolerance Prohibited Materials Audit (\texttt{Rule 3.b.i} \& FAQ 10/22/2026 on polystyrene):} Verify that \textbf{zero Polystyrene / Styrofoam / \texttt{\#6 PS} plastic}, fiberglass, mineral wool, asbestos, or mercury thermometers are present anywhere in the device or lab kit (\texttt{Rule 3.b.iii} bars any device containing prohibited hazardous materials from testing: $\mathrm{TS}=0, \mathrm{PS}=0$!).
  \item \textbf{Electrical \& Water-Spill Separation:} Plug the clamp lamp into a GFCI-protected $\SI{120}{\volt}$ outlet with a drip loop so water spills on the Testing Surface can never reach electrical cords. Never place or remove the water cup while looking directly into the incandescent filament.
\end{itemize}
\end{rulebox}

\alsosee{\sref{sec:00} (Command Center \& Rules), \sref{sec:02} (Sig Figs \& Units), \sref{sec:14} (Material Optical/Thermal Properties), \sref{sec:15}--\sref{sec:16} (2027 Collector Design \& Calibration DOE), \sref{sec:18} (Graphs, Semi-Log OLS \& Residuals), and \sref{app:D} (Printable Calibration Logs).}
